\hypertarget{ux6458ux8981}{%
\section{摘要}\label{ux6458ux8981}}

海底地形地貌匹配导航利用水下声学观测与环境地图之间的对应关系，为惯性与航位推算提供位置或运动约束。该问题的核心不是寻找一个通用匹配器，而是确定海底环境能够提供哪些独立信息，以及如何在地图不完整、观测相关、姿态不确定和匹配存在歧义时保持估计可信。本报告围绕多波束测深、侧扫声呐、惯性测量、多普勒测速与深度观测，综合水下导航经典综述、测深SLAM、概率地图、侧扫SLAM和全局配准研究，建立问题分类、测量模型、算法比较和证据评价框架。

文献表明，已知地图上的地形辅助导航、无先验地图的声呐SLAM、未知初始位置的全局重定位具有不同的估计目标，不能用同一个精度指标替代。多波束提供几何信息，侧扫主要提供具有视角依赖的回波强度和斜距信息；两者具有互补性，但不能直接按激光雷达与相机的测量模型融合。相关匹配适合批量候选搜索，点质量滤波和粒子滤波适合非线性及多峰位置估计，因子图适合跨时间约束与轨迹一致性，概率地图适合表达测深及插值的不确定性。各类方法的有效性均依赖地形信息、数据关联正确性和噪声模型，而非仅依赖算法规模。

报告进一步分析平坦海底、单向坡面、重复地貌、低重叠条带和地图偏差造成的退化，讨论海底高程与反射率联合表示、地图质量参与匹配、多假设跟踪及信息驱动航迹设计。研究结论是：面向长航程水下导航，应把局部相对估计、全局地图关联、概率可信度和失效恢复分层研究，以真实测量几何和独立验证为依据选择方法。本文是基于公开文献的综合研究与理论分析，不包含具体项目的仿真、部署或性能结果。

关键词 海底地形匹配 地形辅助导航 多波束测深 侧扫声呐 粒子滤波 概率地图 水下SLAM

\hypertarget{ux7b2cux4e00ux7ae0-ux7814ux7a76ux5bf9ux8c61ux4e0eux6587ux732eux8303ux56f4}{%
\section{第一章 研究对象与文献范围}\label{ux7b2cux4e00ux7ae0-ux7814ux7a76ux5bf9ux8c61ux4e0eux6587ux732eux8303ux56f4}}

\hypertarget{ux6c34ux4e0bux5b9aux4f4dux95eeux9898ux7684ux57faux672cux7ed3ux6784}{%
\subsection{1.1 水下定位问题的基本结构}\label{ux6c34ux4e0bux5b9aux4f4dux95eeux9898ux7684ux57faux672cux7ed3ux6784}}

水下导航研究需要首先区分自主航行、状态估计和测绘地理参考。自主航行关心车辆能否沿规划航迹安全执行任务；状态估计关心当前位置、速度、姿态及误差状态；测绘地理参考关心每个测深或声呐图像元素在统一坐标系中的位置。导航轨迹平滑不等于绝对定位正确，局部地图没有明显重影也不等于具有准确的地理坐标。Paull等的经典综述为惯性推算、外部声学定位和环境辅助定位提供了基础分类。\href{https://people.csail.mit.edu/lpaull/publications/Paull_JOE_2014.pdf}{［1］}

惯性和多普勒测速可以形成连续的运动估计，但其误差没有天然的地理参考。海底地形导航的意义在于，利用环境中的持久结构建立附加参照，减少完全依赖时间积分的漂移。然而，``地形提供参照''并非意味着任意一次测量都能确定二维位置；参照是否有效取决于不同位置的观测是否可区分。该问题同时涉及测量物理、地貌统计和状态估计。

Carreño等和Melo与Matos分别从地形导航模型与研究进展角度梳理了这一领域。它们说明，地形匹配不仅包括相关搜索，也包括以地形为观测的贝叶斯滤波，以及依靠历史测量形成地图的测深SLAM。\href{https://eprints.soton.ac.uk/153897/1/Carreno_2010_-_A_survey__on_TBN_for_AUVs.pdf}{［2］}\href{https://doi.org/10.1016/j.oceaneng.2017.04.047}{［3］} 因此，研究算法时应明确地图是否已知、初始位置是否有界、是否需要实时输出，以及是否存在重访区域。

\hypertarget{ux56dbux7c7bux5bb9ux6613ux6df7ux6dc6ux7684ux7814ux7a76ux4efbux52a1}{%
\subsection{1.2 四类容易混淆的研究任务}\label{ux56dbux7c7bux5bb9ux6613ux6df7ux6dc6ux7684ux7814ux7a76ux4efbux52a1}}

第一类是局部地形辅助定位。车辆已有较可信的初始位置，在有限区域内将测深与先验地形匹配，估计位置误差。第二类是全局重定位。初始位置未知或原估计失效，需要在大地图中发现候选，并进一步消除相似地貌的歧义。第三类是声呐里程计。利用时间上不同的环境观测估计相对运动，不要求已有地理参考地图。第四类是SLAM。联合估计运动与环境，利用重复观测建立回环或地图一致性。

\begin{longtable}[]{@{}p{0.22\linewidth}p{0.40\linewidth}p{0.31\linewidth}@{}}
\toprule
研究任务 & 主要输入和输出 & 不能替代的评价内容\tabularnewline
\midrule
\endhead
局部地图匹配 & 有界位姿先验与已知地图 输入 输出地图坐标中的位置更新 & 局部成功不能证明未知初始位置可恢复\tabularnewline
全局重定位 & 大范围候选与连续测量 输入 输出候选位置及可信度 & 检索命中不能代替正确位姿与持续跟踪\tabularnewline
声呐里程计 & 连续或重叠历史观测 输入 输出短时相对位姿 & 绝对位置误差不能独立评价相对运动精度\tabularnewline
声呐SLAM & 相对运动与环境关联 输入 输出轨迹和地图 & 内部一致性不能代替地理参考精度\tabularnewline
\bottomrule
\end{longtable}

上述任务之间可以共享特征、子地图或优化工具，但不是同一问题。Barkby等研究的是测深粒子滤波SLAM，而HUGIN地形导航模块强调利用地形位置更新辅助惯性导航。\href{https://doi.org/10.1177/0278364912459666}{［4］}\href{https://www.navlab.net/Publications/The_HUGIN_AUV_Terrain_Navigation_Module.pdf}{［5］} 将前者的建图成果直接作为后者在未知地图上的绝对定位证据，会混淆研究边界。

\hypertarget{ux6587ux732eux7b5bux9009ux539fux5219}{%
\subsection{1.3 文献筛选原则}\label{ux6587ux732eux7b5bux9009ux539fux5219}}

本报告优先采用IEEE Journal of Oceanic Engineering、IEEE Transactions on Robotics、The International Journal of Robotics Research及IEEE Robotics and Automation Letters等来源中的直接相关研究，以作者公开版本和机构资料补充可读性。会议论文用于提供基础方法或工程细节，官方软件和测量指南用于核对模型及术语。期刊类别只构成筛选依据，不能代替对实验设计和测量适用性的判断。

文献核查分为公开全文、公开方法片段、摘要与元数据、官方工程资料四个层级。对于只能获取摘要的研究，仅讨论其公开任务和总体框架，不推断未公开的参数或实验配置。报告中的理论解释与跨论文比较是综合分析，不伪装成某篇论文的原文结论；未采用影响因子排名来断言某个算法绝对最优。

\hypertarget{ux7b2cux4e8cux7ae0-ux6d77ux5e95ux5730ux8c8cux4fe1ux606fux4e0eux5730ux56feux8d28ux91cf}{%
\section{第二章 海底地貌信息与地图质量}\label{ux7b2cux4e8cux7ae0-ux6d77ux5e95ux5730ux8c8cux4fe1ux606fux4e0eux5730ux56feux8d28ux91cf}}

\hypertarget{ux5730ux5f62ux4e0eux5730ux8c8cux4fe1ux606fux7684ux533aux522b}{%
\subsection{2.1 地形与地貌信息的区别}\label{ux5730ux5f62ux4e0eux5730ux8c8cux4fe1ux606fux7684ux533aux522b}}

地形主要描述几何，例如高程、坡度、曲率、脊线与谷线。地貌还包括地形形成过程、形态类别、底质和结构分布。同一高程曲面可以具有不同沉积物和声学反射性质，因此高程匹配和强度匹配不是相同证据。几何较平坦但存在稳定底质边界的区域，可能对侧扫识别有利；几何起伏明显但声学强度随观测角度变化剧烈的区域，则可能更适合测深约束。

海底匹配所需要的不是单纯``粗糙程度''，而是观测在候选位置之间的可区分性。一条长而平直的海底坡面可能具有较大坡度，但沿等深方向没有明显变化；多个大小相似、间距近似的沙波也可能提供很强局部纹理，却在全局上产生重复候选。研究地貌应同时考虑局部可观性和大范围唯一性。

因此，适航地形评价至少包括高程变化幅度、梯度方向分布、曲率稳定性、空间重复度、特征持续性和地图观测来源。对不同尺度的地形做相同评价也可能产生误导：微小沙纹有利于近距离强度对应，却可能无法在粗分辨率先验地图中出现；大型海脊有利于全局检索，但短时运动对其形态影响不大。

\hypertarget{ux5730ux56feux4e0dux662fux65e0ux8befux5deeux771fux503c}{%
\subsection{2.2 地图不是无误差真值}\label{ux5730ux56feux4e0dux662fux65e0ux8befux5deeux771fux503c}}

数字水深模型往往由不同年代、设备、测线和处理方法的数据合成。栅格间距描述存储采样间隔，不能直接等同于有效分辨率、位置精度或高程精度。GEBCO配套的来源类型标识可以区分不同来源；NOAA ETOPO同样是融合地形产品，而不是每个像元都来自独立高密度测深。\href{https://www.gebco.net/data_and_products/gridded_bathymetry_data/}{［6］}\href{https://www.ncei.noaa.gov/products/etopo-global-relief-model}{［7］}

对于地图匹配，至少要识别四类误差：高程随机误差、水平位置误差、垂直基准偏差和由稀疏数据插值造成的结构误差。随机误差可以在某些条件下平均降低，而后面三类通常具有空间相关性，不能通过增加匹配点数自然消除。海图更新或产品勘误也意味着地图版本必须成为定位研究的可追溯条件。\href{https://www.gebco.net/data-products/gridded-bathymetry-data/data-set-errata}{［8］}

Salavasidis等针对粗地图的研究特别有助于理解地图质量与导航的关系：论文讨论了地图畸变和大量插值区域对滤波稳定性的影响，并研究利用更可信及信息更丰富的区域设计航迹。\href{https://eprints.soton.ac.uk/480480/1/joe.pdf}{［9］} 本报告据此认为，地图质量层应当作为观测可信度的输入，而不是仅在下载数据时检查一次。

\hypertarget{ux5730ux7406ux5750ux6807ux548cux5782ux76f4ux57faux51c6}{%
\subsection{2.3 地理坐标和垂直基准}\label{ux5730ux7406ux5750ux6807ux548cux5782ux76f4ux57faux51c6}}

地形匹配可以在经纬度、投影平面、局部切平面或车辆坐标系中计算，但必须保持单位与变换一致。经纬度栅格在东西方向的实际距离随纬度变化，不能直接把角秒当成等间距米制像素。若通过投影重采样生成米制地图，重采样本身会引入平滑及新的相关性。

深度计给出的水深与测深图中的海底高程也未必采用相同垂直参考。压力转深度需要考虑海水密度等因素，潮位和基准面差异可能产生共同偏移。HUGIN公开技术论文明确讨论了压力转深度、潮位及深度偏差问题。\href{https://www.navlab.net/Publications/The_HUGIN_AUV_Terrain_Navigation_Module.pdf}{［5］} 因而匹配中的共同高程偏差必须被校正、估计或消去，不能默认为零。

IHO S-44提供海道测量不确定度与调查质量的规范化语言，可帮助区分测量精度和数据用途。\href{https://iho.int/uploads/user/pubs/standards/s-44/S-44_Edition_6.2.0_adopted.pdf}{［10］} 标准中的概念可以指导误差预算，但不能仅凭使用了某种配准算法就宣称达到某个测量等级。

\hypertarget{ux5730ux56feux8868ux8fbeux7684ux9009ux62e9}{%
\subsection{2.4 地图表达的选择}\label{ux5730ux56feux8868ux8fbeux7684ux9009ux62e9}}

规则栅格具有查询方便、便于多尺度处理和内存可预测的优点，适合高程函数查询与相关搜索。其主要限制是难以自然表示陡壁、悬垂结构和同一水平位置的多个表面，也容易把缺测填充值误当成真实表面。对于大多数海底近似单值地形，栅格仍然是合理基础。

点云和三角网能够更灵活保留测量几何，但匹配时需要处理点密度、局部法向和边界。若地图与实时条带密度不同，最近邻残差可能由采样差异而非实际位姿误差主导。概率表达则可进一步存储高程方差、测量位置不确定性和插值可信度。Torroba等的SVGP地形建模把输入位置的不确定性也纳入地图，具有直接的海底应用依据。\href{https://arxiv.org/abs/2203.10893}{［11］}

地图表达可以多层共存：全局使用低分辨率高程和质量层，局部使用点云或网格，匹配时查询连续概率表面。关键是明确各层来源与相关性，避免把同一原始测深形成的多种表达当作相互独立的证据。

\hypertarget{ux7b2cux4e09ux7ae0-ux6c34ux4e0bux4f20ux611fux5668ux7684ux6d4bux91cfux51e0ux4f55}{%
\section{第三章 水下传感器的测量几何}\label{ux7b2cux4e09ux7ae0-ux6c34ux4e0bux4f20ux611fux5668ux7684ux6d4bux91cfux51e0ux4f55}}

\hypertarget{ux591aux6ce2ux675fux6d4bux6df1ux89c2ux6d4b}{%
\subsection{3.1 多波束测深观测}\label{ux591aux6ce2ux675fux6d4bux6df1ux89c2ux6d4b}}

多波束测深通常在一次声学发射及回波处理周期中形成多个方向上的测距结果。每条测量在声呐坐标中由斜距和波束方向定义，再经安装外参及车辆姿态变换到导航坐标。波束数量、有效开角、测距方式和时间标签由具体设备决定，不能由算法里的子帧长度反推物理发射结构。

若声呐坐标点为p，车辆位姿为T，先验高程函数为H，则可以构造变换后点的垂向坐标与地图高程之间的差。下式采用向上为正的高程约定，b为模型中的共同垂向偏差。采用向下为正的深度约定时，需要一致地修改符号，不能混用。

\begin{figure}
\centering
\includegraphics{assets/map_residual.png}
\caption{式3-1 地形高程残差和局部信息矩阵}
\end{figure}

这个模型看似只有高程残差，却能通过地图梯度约束水平位置：位置向某方向变化会引起预测高程变化。该作用完全依赖局部坡面方向。如果所有有效测量都落在梯度近似相同的区域，某些水平运动方向仍然缺乏信息。多波束数目增加不必然增加信息矩阵的秩。

外侧波束的误差通常受入射角、声速剖面、波束足迹和底面坡度影响。将所有测量赋予同一个固定方差，会隐藏真实的异方差。概率多波束配准研究关注测点及位姿不确定性的传播，为这类问题提供了方法依据。\href{https://pmc.ncbi.nlm.nih.gov/articles/PMC4851074/}{［12］}

\hypertarget{ux4fa7ux626bux58f0ux5450ux89c2ux6d4b}{%
\subsection{3.2 侧扫声呐观测}\label{ux4fa7ux626bux58f0ux5450ux89c2ux6d4b}}

侧扫声呐的图像通常沿一维方向记录回波时间或斜距，另一维由车辆前进过程中的多个ping组成。它不是由同一时刻的二维透视投影形成。某个像素强度来自波束覆盖和回波积分，受声学反射率、入射角、传播损失、增益、阴影和散斑等共同影响。

GeoHab专业指南强调回波强度处理中的设备响应、入射角与校正等问题。\href{https://geohab.org/backscatter-working-group/}{［13］} 这意味着真实强度的稳定性需要辐射校正和观测几何支持。把高程坡度直接转换成一幅灰度图，再假设它等价于真实侧扫反射率，不足以支撑声学匹配研究。

侧扫单像素通常不能唯一确定三维散射点，因为缺少完整的俯仰或高程信息。斜距校正需要离底高度及地形假设；平底近似在斜坡上可能产生系统性地理定位误差。稀疏和稠密侧扫SLAM论文采用声呐特定几何约束，而不是直接套用相机重投影。\href{https://arxiv.org/abs/2304.01854}{［14］}\href{https://arxiv.org/abs/2312.13802}{［15］}

\hypertarget{ux591aux6ce2ux675fux4e0eux4fa7ux626bux7684ux4e92ux8865ux6761ux4ef6}{%
\subsection{3.3 多波束与侧扫的互补条件}\label{ux591aux6ce2ux675fux4e0eux4fa7ux626bux7684ux4e92ux8865ux6761ux4ef6}}

多波束可以为侧扫观测提供地形与高程先验，使侧扫对应关系更容易解释为相对位姿约束。侧扫的底质边界、结构纹理和阴影模式可能帮助区分几何相似的地点。因此，``几何提供结构，强度提供辨识''是有研究价值的融合思路。

但同一声学发射来源不代表两个输出的每个采样位置完全相同，也不保证噪声独立。测深波束中心、回波积分区间、左右舷覆盖、频率与波束足迹均影响对应关系。研究应使用设备的时间标签、波束定义和输出格式建立关联，而不是按图像尺寸逐像素绑定。

多模态融合还要处理共享误差。姿态误差或声速偏差可能同时影响几何点与强度地理参考，若将两者作为完全独立观测，会过度增加信息。声呐学习综述指出的域差异与鲁棒性问题也适用于融合数据：模型在一种设备或处理链路上有效，不等于在另一种设备上具有相同辨识性。\href{https://arxiv.org/abs/2412.11840}{［16］}

\hypertarget{imu-dvlux548cux6df1ux5ea6ux8ba1ux7684ux4f5cux7528}{%
\subsection{3.4 IMU DVL和深度计的作用}\label{imu-dvlux548cux6df1ux5ea6ux8ba1ux7684ux4f5cux7528}}

IMU提供角速度和比力，支持连续姿态与运动传播；DVL提供相对底面或相对水体的速度信息；深度计约束垂直位置。它们的组合为地形匹配提供姿态、速度和短时运动先验，也承担声呐条带拼接的基础。然而先验参与拼接不等于拼接后的声呐已经提供独立运动信息。

DVL底跟踪和水跟踪必须分开建模。底跟踪可用于地速约束，水跟踪与海流相关，若直接当作地速，可能产生持续位置偏差。DVL安装误差、杆臂和速度比例问题也需要考虑。Pan与Wu的自标定研究说明，参数是否能被估计与运动激励和可观性有关。\href{https://arxiv.org/abs/1509.02054}{［17］}

深度计可以显著降低垂直方向不确定性，却不能替代航向约束。重力参考通常有助于滚转和俯仰估计，但不自动决定水平航向。由此可见，水下融合研究需要分析每种传感器真正约束的状态方向，而不是仅统计传感器种类。

\hypertarget{ux7b2cux56dbux7ae0-ux5730ux5f62ux4fe1ux606fux53efux89c2ux6027ux4e0eux9000ux5316ux5206ux6790}{%
\section{第四章 地形信息可观性与退化分析}\label{ux7b2cux56dbux7ae0-ux5730ux5f62ux4fe1ux606fux53efux89c2ux6027ux4e0eux9000ux5316ux5206ux6790}}

\hypertarget{ux5e73ux5766ux6d77ux5e95ux4e0eux5355ux5411ux5761ux9762ux7684ux6839ux672cux9650ux5236}{%
\subsection{4.1 平坦海底与单向坡面的根本限制}\label{ux5e73ux5766ux6d77ux5e95ux4e0eux5355ux5411ux5761ux9762ux7684ux6839ux672cux9650ux5236}}

对水平位置做局部线性化，地形残差的雅可比由地图梯度决定。平坦海底梯度接近零，位置变化几乎不改变预测高程，测深匹配很难修正水平误差。单向坡面虽然梯度不为零，但所有梯度方向相近，通常只能较强约束坡向，而不能充分约束沿等深线的移动。

这是观测结构限制，不能通过更小的迭代步长、更高阶优化器或无限增加同类测点彻底解决。在数值优化中，退化表现为信息矩阵小特征值；在粒子滤波中，表现为沿某方向的权重分布不明显收缩。不同计算框架只是以不同形式呈现同一个信息不足问题。

\hypertarget{ux822aux5411ux53efux89c2ux6027}{%
\subsection{4.2 航向可观性}\label{ux822aux5411ux53efux89c2ux6027}}

航向变化使条带中远离旋转中心的点发生水平位移。若这些位移穿过不同方向的地形梯度，会产生可辨识高程变化；如果地形近似平面、条带过窄或点集中在局部，航向仍可能与平移或垂向偏差耦合。不能因为残差包含旋转变量，就认定旋转已经被充分观测。

转弯同时产生有利和不利作用。不同观察方向可能改善地形信息，但角速度增加会放大时间偏差、杆臂及条带运动畸变。应把``转弯提供激励''和``转弯恶化未校正模型''分开讨论。Xu等关于成像声呐关键帧的工作以雅可比及奇异值分析约束情况，可作为这种结构性分析的参考，但其成像声呐场景不能直接代替测深侧扫场景。\href{https://arxiv.org/abs/2106.16032}{［18］}

\hypertarget{ux5171ux540cux6df1ux5ea6ux504fux5deeux4e0eux6c34ux5e73ux4f4dux7f6eux8026ux5408}{%
\subsection{4.3 共同深度偏差与水平位置耦合}\label{ux5171ux540cux6df1ux5ea6ux504fux5deeux4e0eux6c34ux5e73ux4f4dux7f6eux8026ux5408}}

如果未知垂向偏差b与某方向的地形高度变化相似，那么匹配可以通过改变水平位置来解释垂向误差。估计中加入b会更符合实际，但可能降低位置的信息量；固定b为零可以产生较强数值约束，却可能把错误基准转化为水平偏差。

\begin{figure}
\centering
\includegraphics{assets/nuisance.png}
\caption{式4-1 对共同垂向偏差边缘化后的有效位置信息}
\end{figure}

该式是线性化联合估计中的Schur补关系。它说明可用于位置的净信息需要扣除与干扰参数共享的部分。在实际研究中，可以对比原始高程匹配、去均值轮廓匹配和联合偏差估计，但需要同时解释各自损失了什么信息。去均值可以减少共同偏差，却也删除绝对高程这一有用识别特征。

\hypertarget{ux91cdux590dux5730ux8c8cux4e0eux5168ux5c40ux6b67ux4e49}{%
\subsection{4.4 重复地貌与全局歧义}\label{ux91cdux590dux5730ux8c8cux4e0eux5168ux5c40ux6b67ux4e49}}

局部可观性良好不保证全局唯一性。两个相隔很远的山谷可以各自形成尖锐的匹配峰，局部优化在任意一个峰附近都可能给出很小的协方差。这个协方差只描述所选峰附近的不确定性，并没有表达选择错误地点的风险。

Global UTAN的正式全文说明了如何以VLAD检索先验子图、用TEASER++配准生成候选，并由cohort粒子滤波管理多个位置假设。\href{https://doi.org/10.1109/TIM.2026.3716492}{［19］} 它先依据前\(m\)个检索子图是否覆盖当前估计来判定KRP，再利用运动预测距离与测深MSE选择跟踪输出、抑制无效cohort；因此确实把地点歧义纳入持续跟踪，而不是一次匹配后立即宣布成功。这里必须注意，原文要求已有海底地形先验图，并假设深度、横滚和俯仰误差有界；真实测深数据上的DR噪声和KRP为离线模拟注入，不能等同于AUV艇载闭环验证。

\hypertarget{ux7b2cux4e94ux7ae0-ux76f8ux5173ux5339ux914dux4e0eux5c40ux90e8ux914dux51c6}{%
\section{第五章 相关匹配与局部配准}\label{ux7b2cux4e94ux7ae0-ux76f8ux5173ux5339ux914dux4e0eux5c40ux90e8ux914dux51c6}}

\hypertarget{ux8f6eux5ed3ux548cux533aux57dfux76f8ux5173ux641cux7d22}{%
\subsection{5.1 轮廓和区域相关搜索}\label{ux8f6eux5ed3ux548cux533aux57dfux76f8ux5173ux641cux7d22}}

相关匹配在候选位置、航向或高程偏差范围内，对实时观测与地图预测进行比较。常见目标包括平方误差、绝对误差和归一化相关。平方误差与高斯误差模型相联系，绝对误差相对不易被少量异常值主导，归一化相关更关注形状而弱化幅值，但可能丢失绝对高程信息。

相关搜索的优点是直观、可以展示多个峰，并可构建粗到细搜索。其缺点是搜索维度增加后计算量迅速上升，而且采样步长、地图插值和观测尺度都会影响峰值位置。在地图覆盖不足或真实位置不在搜索范围内时，即使优化返回一个最小值，也不代表存在正确匹配。

地形相关研究应保存完整候选分布或至少保存多个局部最优，而不只保存最高分。第一峰和第二峰的差距只有在候选生成及评分模型一致时才有意义。不同有效测量数量、不同覆盖比例或不同噪声方差的候选，不能直接比较未经规范化的残差和。

\hypertarget{icpux4e0egicp}{%
\subsection{5.2 ICP与GICP}\label{icpux4e0egicp}}

ICP在当前变换下寻找对应点，再估计变换并重复。它适合已有较好初值和足够重叠的局部配准，不应被当作未知位置情况下的全局搜索器。平面海底中，最近邻距离降低并不一定意味着沿平面方向的位置被确定。

GICP以局部协方差描述点云表面结构，把点到点与局部平面关系纳入概率配准框架。\href{https://www.robots.ox.ac.uk/~avsegal/resources/papers/Generalized_ICP.pdf}{［20］} 在海底测深中，该思想有助于区分强约束与弱约束方向，但协方差还需反映实际测量误差和点密度，而不能仅表示几何平面性。

子地图之间的配准信息应来自真正重叠的环境。只有相邻但不重叠的条带，无法通过普通点云对应产生可靠的同点约束。高斯过程外推等方法可以形成额外预测，但其可信度来自空间相关模型及预测方差，不是凭空创造共同测点。\href{https://www-personal.acfr.usyd.edu.au/sbarkby/2011_iros_bpslam.pdf}{［21］}

\hypertarget{ux9c81ux68d2ux5168ux5c40ux914dux51c6}{%
\subsection{5.3 鲁棒全局配准}\label{ux9c81ux68d2ux5168ux5c40ux914dux51c6}}

全局配准需要解决较大初始位姿误差和高比例错误对应。TEASER采用截断最小二乘及可认证估计思想，是重要的鲁棒配准方法。\href{https://arxiv.org/abs/2001.07715}{［22］} 其作用是根据给定对应关系估计变换，不负责保证每个候选地点正确；可认证性也具有模型、噪声边界和求解条件，不能扩大解释为对任意海底数据均能认证正确位置。

分支定界类方法在确定的搜索范围和评分函数内剪枝搜索，可用于不依赖良好初值的全局匹配。3D-BBS给出了点云全局定位的具体实现参考。\href{https://arxiv.org/abs/2310.10023}{［23］} 这类方法的数学最优性与真实地理正确性不是同一件事：若地图有偏差或评分偏爱错误地貌，即使找到了评分意义上的最优，也可能定位错误。

\hypertarget{ux7b2cux516dux7ae0-ux975eux9ad8ux65afux5730ux5f62ux5bfcux822aux4e0eux6982ux7387ux5730ux56fe}{%
\section{第六章 非高斯地形导航与概率地图}\label{ux7b2cux516dux7ae0-ux975eux9ad8ux65afux5730ux5f62ux5bfcux822aux4e0eux6982ux7387ux5730ux56fe}}

\hypertarget{ux8d1dux53f6ux65afux4f30ux8ba1ux7684ux7edfux4e00ux5f62ux5f0f}{%
\subsection{6.1 贝叶斯估计的统一形式}\label{ux8d1dux53f6ux65afux4f30ux8ba1ux7684ux7edfux4e00ux5f62ux5f0f}}

地图匹配可以作为观测似然，惯性与里程计作为状态转移。贝叶斯递推先传播位置分布，再用地形观测改变候选权重。以下关系给出一般形式，并不预设只能使用某一种滤波器。

\begin{figure}
\centering
\includegraphics{assets/bayes.png}
\caption{式6-1 已知地图条件下的地形导航贝叶斯递推}
\end{figure}

卡尔曼类滤波以单峰近似表达状态，适合初始误差较小及局部可线性化的情况。点质量滤波在网格上表达概率，粒子滤波用带权样本近似分布，能够保留相隔较远的多个地点候选。选型时首先要看分布形状及状态维度，再比较实现效率。

\hypertarget{ux70b9ux8d28ux91cfux6ee4ux6ce2}{%
\subsection{6.2 点质量滤波}\label{ux70b9ux8d28ux91cfux6ee4ux6ce2}}

点质量滤波对有限离散位置网格递推概率。它可以直观保留多峰分布，覆盖充分时不易因随机采样遗漏较小候选。Ånonsen与Hallingstad比较了点质量滤波和粒子滤波的水下地形导航应用，提供了直接方法依据。\href{https://www.navlab.net/Publications/Terrain_Aided_Underwater_Navigation_Using_Point_Mass_And_Particle_Filters.pdf}{［24］}

网格方式的主要限制是状态维度。二维位置尚可建立适当分辨率网格，一旦同时加入航向、高程偏差和更多系统误差，网格数量会乘法增长。粗到细、可变区域和分层状态估计有助于降低复杂度，但也可能在错误早期选择后遗漏真实峰。

\hypertarget{ux7c92ux5b50ux6ee4ux6ce2ux53caux5176ux5931ux6548ux65b9ux5f0f}{%
\subsection{6.3 粒子滤波及其失效方式}\label{ux7c92ux5b50ux6ee4ux6ce2ux53caux5176ux5931ux6548ux65b9ux5f0f}}

粒子滤波允许非线性运动和非高斯观测，但有限粒子数会带来采样误差。重要性分布、权重归一化、重采样和粒子多样性共同决定效果。Arulampalam等的教程为这些过程提供理论基础，Fox的KLD采样研究则讨论依据分布复杂度调整粒子数。\href{https://doi.org/10.1109/78.978374}{［25］}\href{https://papers.neurips.cc/paper_files/paper/2001/hash/c5b2cebf15b205503560c4e8e6d1ea78-Abstract.html}{［26］}

在重复地貌中，强而错误的似然会迅速提高某个错误候选的权重，随后重采样可能删除真实候选。增加粒子只能部分降低遗漏概率，不能纠正失真的似然。鲁棒Monte Carlo定位关于运动与观测提议分布的研究提示，应当保留重新发现地点的能力，而不是在错误峰内无限精化。\href{https://www.borg.cc.gatech.edu/projects/monte-carlo-localization.html}{［27］}

地形粒子滤波的权重还可能受到波束相关性影响。一次ping中若所有点受到同一个深度偏差影响，把这些点当作独立观测会把共同误差重复计入，造成极度尖锐的权重。研究应使用联合协方差、分组似然或有效独立信息量，而不能仅把概率连乘次数当成证据数量。

\begin{figure}
\centering
\includegraphics{assets/corr.png}
\caption{式6-2 共同偏差相关模型和高斯联合负对数似然}
\end{figure}

上式的第一项反映残差大小，第二项反映不同协方差对应的概率密度规范化。比较具有不同预测不确定度的候选时，只用残差除以方差而忽略行列式项，可能错误偏爱高不确定度区域。该问题属于概率模型一致性，不是经验门限大小的问题。

\hypertarget{ux9ad8ux65afux8fc7ux7a0bux4e0eux6d4bux6df1slam}{%
\subsection{6.4 高斯过程与测深SLAM}\label{ux9ad8ux65afux8fc7ux7a0bux4e0eux6d4bux6df1slam}}

高斯过程通过核函数表达空间相关，对未观测位置给出高程预测和方差。其优势是把地图表示和置信度联系起来，而不是把插值表面当成确定真值。常规GP的预测形式如下，K为训练位置间的协方差，Σ为观测噪声，k星为查询与训练位置之间的协方差。

\begin{figure}
\centering
\includegraphics{assets/gp.png}
\caption{式6-3 常规高斯过程高程预测均值与方差}
\end{figure}

Barkby等把粒子轨迹表示与测深观测日志结合，以降低每个粒子存储独立地图的开销，并利用GP提供匹配预测。\href{https://doi.org/10.1177/0278364912459666}{［4］} 这一类研究的贡献在于处理地图与轨迹的不确定性耦合，不能简单概括成``没有重叠也能精确定位''。外推的有效范围受核函数、训练覆盖和地形连续性限制。

SVGP方法使用诱导变量和变分学习改善大规模建图的可扩展性，并研究测量位置不确定性。\href{https://arxiv.org/abs/2203.10893}{［11］} 需要注意，预测方差表达的是模型假设下的不确定性；当核函数不适合断崖、底面突变或地图跨域时，模型可能给出偏差且不够保守的预测。概率地图也需要独立校准。

\hypertarget{ux7b2cux4e03ux7ae0-ux56e0ux5b50ux56fe-ux4fa7ux626bslamux4e0eux591aux6a21ux6001ux7814ux7a76}{%
\section{第七章 因子图 侧扫SLAM与多模态研究}\label{ux7b2cux4e03ux7ae0-ux56e0ux5b50ux56fe-ux4fa7ux626bslamux4e0eux591aux6a21ux6001ux7814ux7a76}}

\hypertarget{ux56e0ux5b50ux56feux7684ux4f18ux52bfux548cux524dux63d0}{%
\subsection{7.1 因子图的优势和前提}\label{ux56e0ux5b50ux56feux7684ux4f18ux52bfux548cux524dux63d0}}

因子图把状态变量与观测关系分开表示，可将相对运动、地形位置、回环和传感器约束统一为跨时间优化问题。iSAM2通过Bayes树实现增量平滑和映射，是重要的后端基础。\href{https://www.borg.cc.gatech.edu/papers/kaess12ijrr.html}{［28］} 但增量优化器不会自动判断某条海底匹配是不是来自错误位置，数据关联和观测模型仍需上游提供可信证据。

单峰高斯地形约束适合已经确认的候选。一旦多个地点同时合理，把其中一个候选当成非常确定的绝对位置因子，会迫使图优化变形轨迹来满足错误关联。混合分布和无效假设的研究为错误回环提供了更丰富的表达。\href{https://april.eecs.umich.edu/media/pdfs/olson2012rss.pdf}{［29］} 将这些思想迁移到海底绝对匹配时，需要重新定义候选证据及其相关性。

\hypertarget{ux4fa7ux626bux7a00ux758fux4e0eux7a20ux5bc6ux65b9ux6cd5}{%
\subsection{7.2 侧扫稀疏与稠密方法}\label{ux4fa7ux626bux7a00ux758fux4e0eux7a20ux5bc6ux65b9ux6cd5}}

稀疏方法检测显著点或结构边界，通过跨图像对应形成位姿约束。它便于解释匹配点，但在均匀底质、散斑或大面积阴影区域可能没有足够特征。全自动侧扫SLAM框架以真实Hugin数据和声呐特定几何进行研究，为自动匹配提供直接证据。\href{https://arxiv.org/abs/2304.01854}{［14］}

稠密子帧方法从重叠图像区域获取更多对应，减少只依赖少数显著点的问题。2025年JOE侧扫研究进一步讨论子帧匹配、相对位姿和轨迹优化。\href{https://arxiv.org/abs/2312.13802}{［15］} 稠密匹配仍然需要有效重叠、几何一致和可靠强度，不应被描述为在任何相邻ping之间都能够产生独立完整位姿。

声呐图像的稠密并不意味着统计独立。相邻像素可能受同一波束、同一增益或同一阴影边界影响。把大量高相关对应都赋予独立小方差，会使估计看似精确但对模型偏差敏感。因此密度提升应与空间分布、物理可见性和有效信息分析同步研究。

\hypertarget{ux795eux7ecfux8868ux793aux4e0eux5b9aux4f4dux5efaux56feux8026ux5408}{%
\subsection{7.3 神经表示与定位建图耦合}\label{ux795eux7ecfux8868ux793aux4e0eux5b9aux4f4dux5efaux56feux8026ux5408}}

NeuRSS将侧扫神经渲染形成的高程先验与SLAM位姿优化联系起来，采用交替改进定位和地形估计的框架。\href{https://arxiv.org/abs/2405.05807}{［30］} 它为``强度不是只有纹理用途''提供了研究线索：多视角回波可以参与几何重建，而改进位姿反过来有助于重建。

这种联合问题存在几何、反射率与姿态的耦合。网络可能通过改变地形解释位姿误差，也可能通过改变反射率解释错误的入射角。只有回波拟合变好，不能证明海底几何和车辆轨迹均正确。需要独立测深、交叉测线或外部参考来评价不同变量。

\hypertarget{ux9646ux5730ux76f4ux63a5ux878dux5408ux65b9ux6cd5ux7684ux8fc1ux79fbux8fb9ux754c}{%
\subsection{7.4 陆地直接融合方法的迁移边界}\label{ux9646ux5730ux76f4ux63a5ux878dux5408ux65b9ux6cd5ux7684ux8fc1ux79fbux8fb9ux754c}}

FAST-LIO2以原始激光点到地图约束和迭代滤波实现紧耦合估计；FAST-LIVO2进一步使用统一几何地图关联图像块并融合直接光度误差。\href{https://arxiv.org/abs/2107.06829}{［31］}\href{https://arxiv.org/abs/2408.14035}{［32］} 对水下研究最值得借鉴的是统一空间表示、不同模态在同一状态下更新及短时运动补偿。

不可直接照搬的是测量定义：相机像素满足透视投影与特定光度近似，侧扫像素包含斜距与波束积分；激光扫描点的采样结构与测深条带也不同。把输入名字从LiDAR和camera改成MBES和SSS，不会自动得到正确的水下算法。有关迁移的结论应以重新推导的残差、雅可比和可观性为依据。

\hypertarget{ux7b2cux516bux7ae0-ux4ee3ux8868ux6027ux7814ux7a76ux7684ux6bd4ux8f83}{%
\section{第八章 代表性研究的比较}\label{ux7b2cux516bux7ae0-ux4ee3ux8868ux6027ux7814ux7a76ux7684ux6bd4ux8f83}}

\hypertarget{ux5de5ux7a0bux5730ux5f62ux5bfcux822aux8defux7ebf}{%
\subsection{8.1 工程地形导航路线}\label{ux5de5ux7a0bux5730ux5f62ux5bfcux822aux8defux7ebf}}

HUGIN模块和MBARI TRN体现了先验地形、惯性运动和地形位置更新的工程路线。前者公开论文介绍点质量滤波及海试，后者官方项目资料说明真实海底地形参照导航的应用背景。\href{https://www.navlab.net/Publications/The_HUGIN_AUV_Terrain_Navigation_Module.pdf}{［5］}\href{https://www.mbari.org/project/terrain-relative-navigation/}{［33］} 这类证据的重要性在于系统长期使用和传感器链路，而不是只有离线配准结果。

MB-System公开源代码提供了地形导航的可检查实现，包括相关的点质量与粒子滤波模块。\href{https://github.com/dwcaress/MB-System}{［34］} 开源意味着能够检查接口及模型，不意味着下载后在任意声呐、地图和航迹上都具有相同精度。研究时需要记录版本、输入类型、基准面和质量处理策略。

\hypertarget{ux5730ux70b9ux8bc6ux522bux548cux4e0dux786eux5b9aux6027ux5b66ux4e60ux8defux7ebf}{%
\subsection{8.2 地点识别和不确定性学习路线}\label{ux5730ux70b9ux8bc6ux522bux548cux4e0dux786eux5b9aux6027ux5b66ux4e60ux8defux7ebf}}

测深点云的数据驱动回环检测研究直接面向海底特征稀缺及航位推算漂移的问题，构建学习描述子用于地点关联和粗配准。\href{https://arxiv.org/abs/2209.08578}{［35］} 它说明学习可以参与环境关联，但训练区域和点云采样方式的限制必须保留，不能用训练任务中的回环结果证明跨海域全局定位。

PointNetKL关注GICP配准的协方差估计，而不只是配准平均误差。\href{https://arxiv.org/abs/2003.10931}{［36］} 对导航而言，可信地表达弱约束方向十分重要：一个位姿值较好的匹配，如果协方差严重过小，仍可能破坏后续融合。学习协方差应检验校准和域外行为，而不只检验回归损失。

PointNetVLAD与MinkLoc3D来自陆地点云地点识别，分别代表全局描述子聚合和稀疏卷积检索思路。\href{https://arxiv.org/abs/1804.03492}{［37］}\href{https://arxiv.org/abs/2011.04530}{［38］} 它们可作为检索结构参考，但海底点云的近单曲面、条带覆盖和特征重复意味着需要水下数据适配与评价。高Recall并不能独立保证绝对姿态正确。

\hypertarget{ux7b97ux6cd5ux6bd4ux8f83ux4e0dux80fdux8131ux79bbux573aux666f}{%
\subsection{8.3 算法比较不能脱离场景}\label{ux7b97ux6cd5ux6bd4ux8f83ux4e0dux80fdux8131ux79bbux573aux666f}}

\begin{longtable}[]{@{}p{0.23\linewidth}p{0.34\linewidth}p{0.36\linewidth}@{}}
\toprule
方法路线 & 有利条件 & 主要限制与研究重点\tabularnewline
\midrule
\endhead
地形相关搜索 & 地图可信 搜索区域有界 测量轮廓可区分 & 多峰和维度增长 需要保留候选而非只取第一峰\tabularnewline
点质量滤波 & 低维状态 有限覆盖 需要完整位置分布 & 网格维数灾难 需要分层和质量感知\tabularnewline
粒子滤波 & 非高斯先验 多模态地形 连续观测 & 粒子贫化 错误似然 需要候选保护与恢复\tabularnewline
ICP和GICP & 重叠足够 初值合理 局部几何可靠 & 平坦或低重叠退化 需要方向性协方差\tabularnewline
鲁棒全局配准 & 对应中保留足够正确结构 & 不能自行解决地点歧义 需要检索与后续验证\tabularnewline
因子图 & 可信跨时间关联与相对运动 & 错误因子能污染轨迹 需要鲁棒及多假设表达\tabularnewline
概率地图 & 测量来源和不确定性可建模 & 核与模型失配 需要预测方差校准\tabularnewline
学习型检索和匹配 & 有代表性训练数据且输入定义稳定 & 跨设备跨地貌域差异 需要独立外推验证\tabularnewline
\bottomrule
\end{longtable}

不存在能够在所有地貌、初始误差和传感器条件下同时占优的方法。研究中所谓``先进''应体现更合理的信息利用、误差表达和失效恢复，而不是单纯由是否使用神经网络、粒子滤波或图优化决定。

\hypertarget{ux7b2cux4e5dux7ae0-ux6587ux732eux8bc1ux636eux7684ux8bc4ux4ef7ux65b9ux6cd5}{%
\section{第九章 文献证据的评价方法}\label{ux7b2cux4e5dux7ae0-ux6587ux732eux8bc1ux636eux7684ux8bc4ux4ef7ux65b9ux6cd5}}

\hypertarget{ux7cbeux5ea6-ux7a33ux5b9aux6027ux4e0eux5b8cux6574ux6027}{%
\subsection{9.1 精度 稳定性与完整性}\label{ux7cbeux5ea6-ux7a33ux5b9aux6027ux4e0eux5b8cux6574ux6027}}

绝对轨迹误差衡量与地理参考真值之间的偏差，相对位姿误差衡量指定时间或距离间隔的运动误差，地图一致性衡量重复测量是否对齐。三者具有不同含义，需要分别报告。仅通过整体刚体对齐后的轨迹误差，不能证明绝对重定位有效；仅通过交叉测线吻合，也不能证明全局地理参考正确。

平均误差会隐藏少量严重失败。完整评价应包括误差分布、较高分位误差、最大误差、持续失锁时间、恢复时间和错误地点确认次数。对于海底全局定位，错误而自信的输出通常比暂时保留多个候选更危险，因此可信度校准和完整性必须参与方法比较。

\hypertarget{ux516cux5e73ux6bd4ux8f83ux7684ux6761ux4ef6}{%
\subsection{9.2 公平比较的条件}\label{ux516cux5e73ux6bd4ux8f83ux7684ux6761ux4ef6}}

不同方法至少应使用一致的测量序列、地图版本、初始不确定范围和外部参考。若一种方法获得精确初值而另一种方法需要全局发现位置，其定位误差不能直接代表算法优劣。若一种方法采用最终离线平滑，另一种采用在线实时输出，同样需要区分。

学习方法还必须按海域、测量航次或设备划分训练与测试，而不只是随机分割同一地图的小块。空间相邻样本共享地貌和误差，随机划分容易高估泛化。关于空间结构交叉验证的研究为这种评价设计提供了方法依据。\href{https://doi.org/10.1111/ecog.02881}{［39］}

\hypertarget{ux771fux5b9eux6570ux636eux4e0eux5408ux6210ux6570ux636eux7684ux4f5cux7528}{%
\subsection{9.3 真实数据与合成数据的作用}\label{ux771fux5b9eux6570ux636eux4e0eux5408ux6210ux6570ux636eux7684ux4f5cux7528}}

合成数据能够精确控制偏差、噪声和故障，用于研究可观性或边界；真实海试能够反映测量链路中的声速变化、条带噪声和域差异。两类证据相互补充，不可相互冒充。声呐学习综述将鲁棒性、数据集和合成到真实差异作为重点问题，这也是评价研究成熟度的重要依据。\href{https://arxiv.org/abs/2412.11840}{［16］}

文献中报告的数值只有连同场景、输入、初值及评价定义才可解释。本报告不跨论文拼接不同数据集上的误差数字形成排行榜，也不把公开摘要中的相对改善转化成对其他系统的绝对精度承诺。

\hypertarget{ux4e13ux9898ux8865ux5145-ux6d77ux5e95ux5730ux8c8cux7279ux5f81ux4e0eux5339ux914dux5c3aux5ea6}{%
\section{专题补充 海底地貌特征与匹配尺度}\label{ux4e13ux9898ux8865ux5145-ux6d77ux5e95ux5730ux8c8cux7279ux5f81ux4e0eux5339ux914dux5c3aux5ea6}}

\hypertarget{ux4eceux5730ux8c8cux540dux79f0ux5230ux5bfcux822aux53efux5229ux7528ux91cf}{%
\subsection{从地貌名称到导航可利用量}\label{ux4eceux5730ux8c8cux540dux79f0ux5230ux5bfcux822aux53efux5229ux7528ux91cf}}

Harris等的全球海底地貌研究及Blue Habitats合作项目，按大尺度环境、起伏类别和具体地貌组织海底信息。NOAA的配套说明进一步区分陆架、陆坡、深海和超深海等基础层，以及峡谷、海山、脊、盆地和台地等特征层。\href{https://bluehabitats.org/}{［40］} 这些类别可以作为研究分层依据，但全球地貌标签的空间尺度通常大于近底声呐窗口，不能直接视为精确定位地标。

地貌分类回答``这一片海底是什么类型''，匹配导航回答``当前位置相对于地图或历史观测在哪里''。前者可以依赖较粗范围的形态和地质背景，后者还需要实际观测足迹内的可辨结构。以下各类导航条件是对测量几何的理论分析，并非全球地貌论文对某种导航性能的实验结论。

对于较平坦的沉积底面，测深高程变化可能较弱；若底质过渡具有稳定声学纹理，侧扫仍可能提供环境关联。对于单向陆坡或长直槽谷，垂直梯度方向与等深方向的信息不对称，应注意横坡约束较强、沿等深方向较弱的情形。对于多向岩脊或峡谷分支，梯度方向更丰富，局部配准潜力较大，但陡壁和阴影会减少实际可见覆盖。

对于海山和独立凹地，完整形态可能较有全球辨识力，而只观测到一个小坡面时，仍可能与其他山体局部相似。对于重复沙波或平行脊，短时配准可能具有清楚的局部峰，却存在一个或多个周期之外的相似解。这说明匹配窗口需要兼顾足够跨度和运动畸变，不能只根据地貌``粗糙''来决定。

\hypertarget{ux53efux6d4bux5730ux5f62ux5c5eux6027ux7684ux5c3aux5ea6ux4f9dux8d56}{%
\subsection{可测地形属性的尺度依赖}\label{ux53efux6d4bux5730ux5f62ux5c5eux6027ux7684ux5c3aux5ea6ux4f9dux8d56}}

高程起伏、坡度、曲率、地形位置指数和表面粗糙度都可以作为描述量。地形位置指数可通过中心高程与指定邻域平均高程的差构造，用于描述局部凸起或凹陷；其含义随邻域半径改变。相同位置在小邻域可以表现为脊顶，在更大邻域又可能处于盆地内部。因此，任何地形属性都应同时报告米制分析尺度，而不只是属性图的像素大小。

坡向分布可以提示观测窗口是否包含多方向坡面，曲率可以提示脊谷变化，方向频谱可以提示重复结构及测线伪影。但这些量只是信息的代理，不能代替实际测量Jacobian。一个地图中的陡壁若被声呐遮挡，对当前定位就没有实际约束；高程来源不可靠的``丰富结构''也不应获得高权重。

对形态描述子，应考虑对旋转、平移和高程偏差的不变性如何影响定位。旋转不变有利于全球检索，却可能丢掉航向信息；去均值有利于消除共同深度偏差，却降低绝对水深的辨识力。检索描述子可以主动去除某些干扰，而精化阶段需要恢复与状态估计有关的物理量，两阶段不必使用完全相同表示。

\hypertarget{ux5730ux8c8cux5206ux5c42ux8bc4ux4ef7ux7684ux7814ux7a76ux610fux4e49}{%
\subsection{地貌分层评价的研究意义}\label{ux5730ux8c8cux5206ux5c42ux8bc4ux4ef7ux7684ux7814ux7a76ux610fux4e49}}

把所有海底区域合并为一个平均误差，容易由信息丰富区域掩盖平坦和重复区域的失败。研究应至少按有效起伏、梯度方向多样性、重复性和观测覆盖分层，并单独记录由地图缺测或来源不可靠引起的失效。地貌类别可帮助解释环境，但不应把类别和导航困难程度建立无条件的一一映射。

多模态研究尤其需要保留``几何丰富但纹理弱''``几何弱但纹理丰富''``两者均丰富''和``两者均弱''的条件区分。只有在这些层次上比较几何、强度及联合方法，才能识别互补发生在哪里。声学底质反射和辐射校正可依据GeoHab专业指南核对，具体信息互补仍需任务级验证。\href{https://geohab.org/backscatter-working-group/}{［13］}

\hypertarget{ux4e13ux9898ux7814ux7a76ux4e00-ux5730ux5f62ux7edfux8ba1ux4e0eux5b9aux4f4dux4fe1ux606fux7684ux63a8ux5bfc}{%
\section{专题研究一 地形统计与定位信息的推导}\label{ux4e13ux9898ux7814ux7a76ux4e00-ux5730ux5f62ux7edfux8ba1ux4e0eux5b9aux4f4dux4fe1ux606fux7684ux63a8ux5bfc}}

\hypertarget{ux5730ux5f62ux7c97ux7cd9ux5ea6ux4e0dux80fdux4ee3ux66ffux4fe1ux606fux77e9ux9635}{%
\subsection{地形粗糙度不能代替信息矩阵}\label{ux5730ux5f62ux7c97ux7cd9ux5ea6ux4e0dux80fdux4ee3ux66ffux4fe1ux606fux77e9ux9635}}

高程方差只描述区域起伏幅度，不能描述不同位移如何影响观测。假设海底在局部近似为线性坡面，高程沿一个方向线性变化，沿另一个方向保持不变。即使坡度很大，沿不变方向的所有平移仍具有相同高程预测。因此该区域的粗糙度可以较高，但二维位置并未完整可观。

若观测覆盖包含多个不同坡向，则梯度外积的累积可能在两个水平维度形成信息。由此可解释为什么具有交叉脊谷的区域通常比单向平坡更有利，也可以解释为什么增加相同方向的测深点主要增加已有方向的信息强度，而不是增加新的独立方向。这个推导是通用残差线性化的结论，并非某篇论文独有的经验规则。

分析信息矩阵时，最小特征值可以提示弱方向，条件数可以提示方向失衡，但必须考虑状态单位和测量协方差。位置与角度混合的矩阵未经尺度规范化，其条件数很难直接比较。对位置、姿态和垂向偏差进行分块分析，比仅报告一个总条件数更清楚。

\hypertarget{ux5c40ux90e8ux4fe1ux606fux4e0eux5730ux56feux7684ux5168ux5c40ux552fux4e00ux6027}{%
\subsection{局部信息与地图的全局唯一性}\label{ux5c40ux90e8ux4fe1ux606fux4e0eux5730ux56feux7684ux5168ux5c40ux552fux4e00ux6027}}

局部信息矩阵只描述某个候选附近的曲率。两个完全相似的海底凸起可以各自提供良好曲率，但无法在全局上判断当前是哪一个。局部可观性和全局辨识性因此应使用不同的量：前者分析雅可比和误差方向，后者分析不同地图区域的描述子或序列似然。

一个区域在局部能产生很多高质量特征，却在整个地图中多次重复，这种区域适合相对配准但可能不适合单帧绝对初始化。反过来，一个大型缓慢变化的海盆深度分布可能有利于粗全局识别，但短时位移使观测变化很小，不利于里程计。这一对比可以解释为何全局候选和局部前端需要不同尺度表示。

地点识别研究使用学习描述子寻找历史地形，对局部配准提供候选；它并不自动替代运动估计。\href{https://arxiv.org/abs/2209.08578}{［35］} 研究应把检索召回、位姿精化和连续确认分别评价，从而避免把一个环节的成功扩展为全部定位链路的成功。

\hypertarget{ux5730ux5f62ux9891ux7387ux4e0eux4f20ux611fux5668ux8db3ux8ff9}{%
\subsection{地形频率与传感器足迹}\label{ux5730ux5f62ux9891ux7387ux4e0eux4f20ux611fux5668ux8db3ux8ff9}}

海底结构具有多个空间尺度。大尺度低频起伏较容易在粗地图中保留，小尺度地貌可能只存在于高密度近底测深或侧扫图中。传感器足迹和制图平滑相当于对部分高频进行抑制，若实时观测与地图保留的频段不同，直接匹配会产生结构不一致。

多尺度匹配的合理目的，是先使用双方共同具有的尺度建立候选，再在可靠细节上精化，而非一味提高图像分辨率。粗地图研究显示，地图畸变和插值质量对导航有实际影响。\href{https://eprints.soton.ac.uk/480480/1/joe.pdf}{［9］} 在没有独立观测支持时，超分生成的微小结构不应自动成为强位置参照。

\hypertarget{ux4e13ux9898ux7814ux7a76ux4e8c-ux6982ux7387ux8befux5deeux4e0eux6709ux754cux96c6ux5408ux4f30ux8ba1}{%
\section{专题研究二 概率误差与有界集合估计}\label{ux4e13ux9898ux7814ux7a76ux4e8c-ux6982ux7387ux8befux5deeux4e0eux6709ux754cux96c6ux5408ux4f30ux8ba1}}

\hypertarget{ux6982ux7387ux63cfux8ff0ux7684ux9002ux7528ux5047ux8bbe}{%
\subsection{概率描述的适用假设}\label{ux6982ux7387ux63cfux8ff0ux7684ux9002ux7528ux5047ux8bbe}}

概率滤波需要明确误差分布。高斯模型便于递推，但海底地图的系统偏差和局部缺测不一定是独立高斯量。若已知误差主要具有可信边界，而难以获得稳定概率分布，集合估计可以提供不同的研究路线。

Ma等在Journal of Field Robotics发表的粗地图多zonotope地形导航研究，以集合成员滤波处理粗地图及自相似地形。\href{https://onlinelibrary.wiley.com/doi/abs/10.1002/rob.22407}{［41］} 其公开摘要说明了问题和总体方法，本报告不推断未公开算法细节。该文值得纳入调研，因为它说明粒子概率不是表达地图歧义的唯一方式。

\hypertarget{ux6982ux7387ux533aux95f4ux4e0eux8befux5deeux8fb9ux754cux7684ux533aux522b}{%
\subsection{概率区间与误差边界的区别}\label{ux6982ux7387ux533aux95f4ux4e0eux8befux5deeux8fb9ux754cux7684ux533aux522b}}

概率置信区间表达在某一统计模型下的覆盖程度，有界集合表达满足误差边界的可能状态。二者不可互换。概率模型失配会造成实际覆盖率低于宣称值；集合方法的边界若漏掉真实误差，则可能排除真实位置。

集合方式也存在保守性问题。误差边界过大时，候选集合可能难以收缩，实际定位不可用；过小则容易虚假收敛。多集合表示可保留分离候选，但数量管理和更新近似同样需要研究。方法选型应依据数据可提供何种误差证据，而不是把某种数学形式视为天然更可靠。

\hypertarget{ux8befux5deeux6765ux6e90ux7684ux5206ux5c42}{%
\subsection{误差来源的分层}\label{ux8befux5deeux6765ux6e90ux7684ux5206ux5c42}}

测距随机误差、单ping共同偏差、测线系统偏差和地图产品基准偏差具有不同时间及空间尺度。随机误差可以在独立条件下累计，持续共同偏差需要状态或相关模型，地图产品偏差需要地图层描述。把所有误差合并为一个数值，会隐藏它们对匹配和恢复的不同影响。

HUGIN公开论文讨论了地形位置更新中可能存在的相关偏差，专业测量标准则规范了测量不确定度语言。\href{https://www.navlab.net/Publications/The_HUGIN_AUV_Terrain_Navigation_Module.pdf}{［5］}\href{https://iho.int/uploads/user/pubs/standards/s-44/S-44_Edition_6.2.0_adopted.pdf}{［10］} 研究中可以以这些资料为基础建立误差预算，但仍需根据实际观测条件确定数值，不应仅复制文献中的参数。

\hypertarget{ux4e13ux9898ux7814ux7a76ux4e09-ux5730ux56feux8868ux793aux4e0eux5728ux7ebfux5b66ux4e60ux7684ux7814ux7a76ux6f14ux8fdb}{%
\section{专题研究三 地图表示与在线学习的研究演进}\label{ux4e13ux9898ux7814ux7a76ux4e09-ux5730ux56feux8868ux793aux4e0eux5728ux7ebfux5b66ux4e60ux7684ux7814ux7a76ux6f14ux8fdb}}

\hypertarget{ux4eceux8f68ux8ff9ux5730ux56feux5230ux5728ux7ebfux6982ux7387ux5730ux56fe}{%
\subsection{从轨迹地图到在线概率地图}\label{ux4eceux8f68ux8ff9ux5730ux56feux5230ux5728ux7ebfux6982ux7387ux5730ux56fe}}

传统每粒子一张完整地图会带来大量重复存储。轨迹地图思想保存每个候选运动历史，与共同观测日志关联，在需要时重建条件地图。Barkby的研究为这一思路提供基础。\href{https://doi.org/10.1177/0278364912459666}{［4］} 该结构利用原始观测共享，但不同轨迹对应不同几何解释，因此不能只共享一张确定表面而忽略候选差异。

Torroba等2023年RA-L工作进一步研究在线SVGP测深SLAM，结合概率地图与高效粒子轨迹表示，公开作者版本和实现资料。\href{https://arxiv.org/abs/2211.05601}{［42］} 它是连接概率地图理论与在线大范围建图的重要文献。此处只概括其公开总体方法，不将论文运行条件扩展为任意设备的保证。

\hypertarget{ux5730ux56feux66f4ux65b0ux548cux67e5ux8be2ux7684ux7edfux8ba1ux89d2ux8272}{%
\subsection{地图更新和查询的统计角色}\label{ux5730ux56feux66f4ux65b0ux548cux67e5ux8be2ux7684ux7edfux8ba1ux89d2ux8272}}

学习地图通常是压缩原始观测的统计模型。对当前位姿查询地图可以产生残差，但如果地图刚使用了同一条测量更新，再用这条测量去``验证''地图，就会出现信息重复。即使从软件接口看是一次训练和一次定位，统计上仍是同一观测。

因此，在线地图与定位耦合需要明确时间顺序和观测身份。联合估计可以显式表达相关性；分离估计需要合理近似及保守协方差。只看地图越来越平滑，不足以判断定位信息是否增加。

\hypertarget{ux53cdux5c04ux7387ux5730ux56feux4e0eux9ad8ux7a0bux5730ux56feux7684ux5206ux5de5}{%
\subsection{反射率地图与高程地图的分工}\label{ux53cdux5c04ux7387ux5730ux56feux4e0eux9ad8ux7a0bux5730ux56feux7684ux5206ux5de5}}

高程、法向和反射率可以作为同一环境的不同属性，但它们不是同一种物理量。侧扫深度重建和法向学习研究试图利用强度恢复几何，结合多视角或其他高程约束。\href{https://arxiv.org/abs/2206.07810}{［43］}\href{https://arxiv.org/abs/2206.07819}{［44］} 这些研究说明强度含有几何信息，也说明几何与底质耦合需要专门处理。

可微侧扫渲染研究从已知或待估表面预测图像，以优化方式恢复参数。\href{https://arxiv.org/abs/2206.07821}{［45］} 对导航的启示是可以用物理可见性解释强度，而不是只比较灰度。它同时带来更复杂的多解性，反射率和地形均可改变回波，因此需要独立约束。

\hypertarget{ux4e13ux9898ux7814ux7a76ux56db-ux6587ux732eux65b9ux6848ux7684ux6761ux4ef6ux5316ux9009ux62e9}{%
\section{专题研究四 文献方案的条件化选择}\label{ux4e13ux9898ux7814ux7a76ux56db-ux6587ux732eux65b9ux6848ux7684ux6761ux4ef6ux5316ux9009ux62e9}}

\hypertarget{ux6309ux5148ux9a8cux548cux73afux5883ux4fe1ux606fux9009ux62e9}{%
\subsection{按先验和环境信息选择}\label{ux6309ux5148ux9a8cux548cux73afux5883ux4fe1ux606fux9009ux62e9}}

若位置先验较可信且地图精度较高，可以优先研究局部地形更新及可信协方差；若初始范围大而地貌重复，应研究多候选全局发现与持续跟踪；若无先验地理地图但存在重访，应研究相对声呐SLAM；若既无地图又无重叠，则不能期待纯环境配准提供持续绝对定位。

这种条件化分类比按``经典''和``先进''划分更有意义。一个成熟点质量滤波器可能比未经域适配的深度检索更适合某个海图条件，学习型方法也可能在有足够真实数据时提供更高召回。比较需要保持输入与任务一致。

\hypertarget{ux6309ux4fe1ux606fux6765ux6e90ux89e3ux91caux6539ux8fdb}{%
\subsection{按信息来源解释改进}\label{ux6309ux4fe1ux606fux6765ux6e90ux89e3ux91caux6539ux8fdb}}

合理改进可以来自更准确的测量模型、更高有效覆盖、更正确的数据关联、更合理的误差表达或更强的恢复机制。它们需要分别证实。仅改变优化器、增加迭代或添加输出平滑，可能改善数值或显示，却不一定增加真实定位信息。

研究报告据此建议，在每条文献路线中都记录所解决的具体问题、需要的输入、保留的假设和独立验证证据。对于没有充分覆盖的情形，应明确开放问题。这种方法可以把文献综述转化为可检验的科学判断，而不是模型名称列表。

\hypertarget{ux6df1ux5316ux7814ux7a76ux4e00-ux771fux5b9eux6d77ux5e95ux5730ux5f62ux6570ux636eux7684ux8bc1ux636eux5206ux7ea7ux4e0eux9009ux53d6}{%
\section{深化研究一 真实海底地形数据的证据分级与选取}\label{ux6df1ux5316ux7814ux7a76ux4e00-ux771fux5b9eux6d77ux5e95ux5730ux5f62ux6570ux636eux7684ux8bc1ux636eux5206ux7ea7ux4e0eux9009ux53d6}}

\hypertarget{ux4e00ux516cux5f00ux5e95ux56feux6d4bux91cfux6805ux683cux4e0eux539fux59cbux6d4bux7ebfux662fux4e0dux540cux8d44ux6599}{%
\subsection{一、公开底图、测量栅格与原始测线是不同资料}\label{ux4e00ux516cux5f00ux5e95ux56feux6d4bux91cfux6805ux683cux4e0eux539fux59cbux6d4bux7ebfux662fux4e0dux540cux8d44ux6599}}

研究数据不能只以``真实海底数据''作为统一标签。第一层是融合底图，适合区域规划和地貌分层，但内部可能包含卫星反演、插值及不同年代的测量。第二层是具有明确来源的测量栅格，适合地形匹配及高程研究，但通常已丢失逐波束时间和部分测量不确定度。第三层是原始或接近原始的声呐测线，能够支持波束级几何、时间同步和导航误差分析。三层的真实性都可能成立，能回答的问题却不相同。

从栅格采样产生观测，只能用于地图匹配方法的受控研究，不能证明声呐驱动程序、运动去畸变或真实强度关联正确。反之，拥有一条真实测线也不能自动获得独立的绝对定位参考。研究资料需要分别标注``地形由实际测量获得''``声学观测由真实设备获得''``参考轨迹由独立系统获得''，不能用其中一项代替其他两项。

\hypertarget{ux4e8cux56dbux7c7bux53efux64cdux4f5cux7684ux516cux5f00ux6570ux636eux5165ux53e3}{%
\subsection{二、四类可操作的公开数据入口}\label{ux4e8cux56dbux7c7bux53efux64cdux4f5cux7684ux516cux5f00ux6570ux636eux5165ux53e3}}

NOAA的CUDEM提供沿海地形与水深产品，官方介绍包括1/9角秒和1/3角秒产品；这些数值表示网格间距，不表示每个位置都具有相同测量精度。其优势是产品、区域和来源较明确，限制是沿海覆盖不能代表深海环境。\href{https://www.ncei.noaa.gov/products/coastal-elevation-models}{［46］}

GMRT融合经过整理的船载多波束资料和背景地形，官方说明区分多波束贡献区域与全球背景。它适合发现测量覆盖较好的海域，不能因名称包含``多分辨率''就认定整个全球产品具有米级细节。\href{https://gmrt.org/about/index.php}{［47］} EMODnet提供欧洲海域DTM、来源及质量信息，并在选定区域提供更高分辨率产品；区域产品的获取条件和垂直基准应以相应版本说明为准。\href{https://emodnet.ec.europa.eu/en/bathymetry}{［48］}

AusSeabed可通过测量覆盖及航次元数据检索水深产品，在部分资料中还可获取相关后向散射数据。水深与强度是否同航次、是否同设备以及是否开放，需要逐项核查；不能从门户存在``backscatter''分类推断每一份水深图都有配套强度。\href{https://www.ausseabed.gov.au/data/bathymetry}{［49］}

\begin{longtable}[]{@{}p{0.20\linewidth}p{0.36\linewidth}p{0.37\linewidth}@{}}
\toprule
数据入口 & 适合支撑的研究 & 下载后首先排除的误解\tabularnewline
\midrule
\endhead
CUDEM & 特定沿海区域的测量栅格、尺度对比及基准核查 & 小网格不等于小测量误差 沿海不等于深海\tabularnewline
GMRT & 按测量贡献选择海域 区域地貌分层 & 背景融合区不等于直接多波束覆盖区\tabularnewline
EMODnet & 欧洲海域地形与来源质量分层 & 高分产品仅覆盖特定区域 质量字段含义不可混用\tabularnewline
AusSeabed & 按航次追溯水深及可能的后向散射 & 覆盖多边形不等于可直接使用的完整观测序列\tabularnewline
\bottomrule
\end{longtable}

\hypertarget{ux4e09ux5efaux7acbux6700ux5c0fux53efux590dux6838ux6570ux636eux6e05ux5355}{%
\subsection{三、建立最小可复核数据清单}\label{ux4e09ux5efaux7acbux6700ux5c0fux53efux590dux6838ux6570ux636eux6e05ux5355}}

每份数据至少记录产品名称和版本、原始访问链接、下载日期、空间范围、投影及轴顺序、水平与垂直单位、垂直基准、网格间距、无数据值、来源类型、调查时间和许可条件。文件校验值用于确认后续研究使用同一份资料；它证明文件一致，不证明数据正确。

质量清单还应记录数据经过的处理。重投影、插值、填洞和平滑都会影响高频地形，尤其要保留处理前文件。空值不得转成零水深；没有覆盖的区域应由掩膜表达。将某个像元内高程的标准差当成测量误差标准差也不合理：前者可能主要反映真实坡度或地形起伏，后者描述对未知真实高程的认识不确定性，两者不是同一随机量。

对于跨航次研究，要检查参考图和待匹配测线是否存在来源交叉。若先验图本身含有待评价测线，将该测线对地图的匹配称为独立外部验证，会高估关联成功率。建议明确设置来源独立、时间独立和测线独立三种标识，并将未知标识保留为未知，而不是补写为独立。

\hypertarget{ux6df1ux5316ux7814ux7a76ux4e8c-ux6d4bux6df1ux7279ux5f81ux65b9ux6cd5ux4e0eux76f4ux63a5ux914dux51c6ux7684ux6761ux4ef6ux6bd4ux8f83}{%
\section{深化研究二 测深特征方法与直接配准的条件比较}\label{ux6df1ux5316ux7814ux7a76ux4e8c-ux6d4bux6df1ux7279ux5f81ux65b9ux6cd5ux4e0eux76f4ux63a5ux914dux51c6ux7684ux6761ux4ef6ux6bd4ux8f83}}

\hypertarget{ux4e00ttt-slamux63d0ux4f9bux7684ux76f4ux63a5ux76f8ux5173ux65b9ux6cd5ux4f9dux636e}{%
\subsection{一、TTT SLAM提供的直接相关方法依据}\label{ux4e00ttt-slamux63d0ux4f9bux7684ux76f4ux63a5ux76f8ux5173ux65b9ux6cd5ux4f9dux636e}}

Zhang与Kim的TTT SLAM采用地形特征、TEASER++和GNC-TLS构成测深SLAM框架。公开内容描述了利用深度变化选择特征、构造旋转不变的地形描述及进行鲁棒关联，并以深度、滚转和俯仰较准确为背景讨论水平位置及航向不确定性。该文为``测深也可以使用特征''提供了直接相关依据，但未公开核读的细节不能据此补写成完整可复现参数。\href{https://www.sciencedirect.com/science/article/pii/S0029801824001148}{［50］}

本报告据此提出的方法比较不是照抄该文实现。对于近似2.5维的海底，特征可以以高程、局部坡度及相对形状表达；这种表示有利于较大位置偏差下寻找对应。直接几何配准则利用较多观测，适合已有合理初值、可见面重叠充分且不确定性可建模的情况。二者分别解决``对应从哪里来''和``已有关联怎样精化''，可以组成两级流程，不必互相排斥。

\hypertarget{ux4e8cux7279ux5f81ux5750ux6807ux548cux63cfux8ff0ux5e94ux5206ux5f00ux8bbeux8ba1}{%
\subsection{二、特征坐标和描述应分开设计}\label{ux4e8cux7279ux5f81ux5750ux6807ux548cux63cfux8ff0ux5e94ux5206ux5f00ux8bbeux8ba1}}

特征位置可以来自脊谷交汇、显著曲率变化、局部深度变化或稳定底质边界。描述子则应编码该位置周围的结构，使同一地点在不同观测下仍可识别。只依据局部高度排序的描述容易被共同深度偏差影响；只依据高度差虽可减少偏置影响，又可能丢掉区分两个海脊的绝对尺度信息。因此应保存原始度量信息与归一化描述，分别用于检索和后续物理核验。

旋转不变并非无条件更好。如果车辆航向已有可靠先验，完全删除方向信息会使不同朝向的重复地貌更相似。合理做法是区分描述阶段的不变性与匹配阶段的先验：检索允许较宽方向变化，精化阶段使用航向概率限制不合理候选，并保留相同地形在多个朝向上的解释。

\hypertarget{ux4e09ux7a20ux5bc6ux65b9ux6cd5ux548cux7279ux5f81ux65b9ux6cd5ux7684ux5171ux540cux8fb9ux754c}{%
\subsection{三、稠密方法和特征方法的共同边界}\label{ux4e09ux7a20ux5bc6ux65b9ux6cd5ux548cux7279ux5f81ux65b9ux6cd5ux7684ux5171ux540cux8fb9ux754c}}

对均匀平面，增加特征检测器不会产生真实角点；提高点云密度也不会产生平面内的独立约束。对重复沙波，特征匹配可能产生多个规则排列的对应，ICP也可能收敛到相邻波峰。对几乎无重叠的测线，任何对应都可能由形状相似而非同一地点产生。因此应先检查空间覆盖和信息可辨识性，再解释算法名称。

\begin{longtable}[]{@{}p{0.22\linewidth}p{0.34\linewidth}p{0.37\linewidth}@{}}
\toprule
研究条件 & 候选机制 & 需要独立验证的环节\tabularnewline
\midrule
\endhead
初值较准且重叠充分 & 直接几何配准或概率配准 & 残差相关、法向精度及退化方向\tabularnewline
初值较差但地貌有显著结构 & 特征检索与鲁棒初始化 & 特征重复度和物理对应的一致性\tabularnewline
地形重复但连续航迹可用 & 多候选连续跟踪 & 时间相关证据和运动先验质量\tabularnewline
平坦且底质稳定可辨 & 声学强度补充关联 & 强度角度依赖和几何映射误差\tabularnewline
无独立结构或地图失配 & 扩大不确定性 保留失效状态 & 不能以优化收敛代替正确定位\tabularnewline
\bottomrule
\end{longtable}

\hypertarget{ux56dbux9c81ux68d2ux7b97ux6cd5ux5e76ux4e0dux80fdux4feeux590dux9519ux8befux95eeux9898ux5b9aux4e49}{%
\subsection{四、鲁棒算法并不能修复错误问题定义}\label{ux56dbux9c81ux68d2ux7b97ux6cd5ux5e76ux4e0dux80fdux4feeux590dux9519ux8befux95eeux9898ux5b9aux4e49}}

TEASER研究为含外点的点云配准提供鲁棒及可认证方法依据。\href{https://arxiv.org/abs/2001.07715}{［22］} GNC研究通过逐步变化的非凸目标处理外点，属于鲁棒估计方法来源，不能一概理解为所有目标都能得到无条件全局最优。\href{https://arxiv.org/abs/1909.08605}{［51］} 这些方法可以降低部分错误对应的影响，却无法判断两个重复地貌中哪一个是实际位置，也无法把没有真实重叠的两片地形变成同一地点。

因此研究比较应报告对应正确率、初值误差范围、实际重叠率以及成功后的精化误差。仅报告``鲁棒优化完成''或者``代价下降''，不足以证明地形关联正确。认证也只针对规定模型及输入条件；声呐几何误差、地图基准错误或系统性声速误差若不在模型内，不会被认证自动消除。

\hypertarget{ux6df1ux5316ux7814ux7a76ux4e09-ux53efux6267ux884cux7684ux8c03ux7814ux8f6cux7814ux7a76ux6d41ux7a0b}{%
\section{深化研究三 可执行的调研转研究流程}\label{ux6df1ux5316ux7814ux7a76ux4e09-ux53efux6267ux884cux7684ux8c03ux7814ux8f6cux7814ux7a76ux6d41ux7a0b}}

\hypertarget{ux4e00ux5148ux5efaux7acbux95eeux9898ux5408ux540cux518dux9009ux62e9ux7b97ux6cd5}{%
\subsection{一、先建立问题合同再选择算法}\label{ux4e00ux5148ux5efaux7acbux95eeux9898ux5408ux540cux518dux9009ux62e9ux7b97ux6cd5}}

研究入口应给出五项条件：地图是否已知、初始位置是否有界、允许使用哪些传感器、能否重访以及需要什么输出。输出可以是水平位置、完整位姿、短时相对变化或带可信度的多个候选。缺少这一合同，``最好的SLAM算法''没有可比较含义。Cadena等的SLAM综述可作为地图、数据关联与估计架构的总框架参考，但不能代替水下测量模型。\href{https://arxiv.org/abs/1606.05830}{［52］}

在多波束、侧扫、IMU、DVL和深度计的传感器集合下，应先从测量几何推导可约束量：IMU提供短时动态传播，DVL提供有条件的速度，深度计提供垂直信息，多波束提供海底几何，侧扫提供强度与斜距相关结构。没有外部地理参考时，局部建图的平移及部分方向自由度不能通过给算法更大网络自然消除。

\hypertarget{ux4e8cux56dbux9636ux6bb5ux7814ux7a76ux53caux9636ux6bb5ux4ea7ux7269}{%
\subsection{二、四阶段研究及阶段产物}\label{ux4e8cux56dbux9636ux6bb5ux7814ux7a76ux53caux9636ux6bb5ux4ea7ux7269}}

第一阶段完成数据和模型核查，产物是坐标约定、时间约定、误差预算与数据来源表。第二阶段建立可解释基线，至少包含运动先验、有限范围高程相关和概率几何配准，产物是输入输出合同及独立评价脚本。第三阶段增加特征、概率地图或多候选跟踪，每次仅改变一个主要因素，产物是条件化对照与失败案例。第四阶段检查长时间、跨区域及失效恢复，产物是成功域、失败域和可信度校准报告，而非只给出单一平均误差。

这四阶段是基于文献综合提出的后续研究步骤，不是已经完成的实验记录。依赖关系明确：没有测量标定不能解释配准残差，没有来源独立的评价不能解释学习增益，没有退化条件不能解释失败是否属于算法缺陷。

\hypertarget{ux4e09ux7ed3ux679cux89e3ux91caux7684ux6700ux5c0fux8868ux683c}{%
\subsection{三、结果解释的最小表格}\label{ux4e09ux7ed3ux679cux89e3ux91caux7684ux6700ux5c0fux8868ux683c}}

建议每个方法报告地图质量类别、地形类别、初始误差、观测覆盖、关联成功率、条件定位误差、失效持续时间和置信区间。总体精度应与定位可用率一同呈现。若某方法只在明显海脊上输出，另一个方法在全航段输出，将两者已输出样本的平均误差直接比较是不公平的。

还应保存无效输出及失败原因。对地形无信息的区间，正确扩大不确定性比给出平滑而错误的高置信度位置更有研究价值。科学结论应表述为``在何种地图、观测和先验条件下具有优势''，而不是声称任何平坦、无重叠或未知地图条件下均可高精度工作。

\hypertarget{ux56dbux8c03ux7814ux540eux7684ux4f18ux5148ux7814ux7a76ux95eeux9898}{%
\subsection{四、调研后的优先研究问题}\label{ux56dbux8c03ux7814ux540eux7684ux4f18ux5148ux7814ux7a76ux95eeux9898}}

优先问题一是条带几何在沿航向和航向角上的信息量怎样随高度、开角和航迹变化。优先问题二是真实强度对几何弱方向是否提供独立信息，以及这种增益在不同底质和入射角下能否保持。优先问题三是地图质量能否通过概率或有界模型进入匹配，使系统在错误地图区不产生虚假收敛。优先问题四是多候选全局重定位与连续相对运动怎样共同工作，而不重复计算同一观测证据。

以上问题将地貌调查、传感器模型、算法研究和评价证据连接起来，也为超分、重定位和声呐里程计三个专题提供统一边界。它们既能分阶段执行，也允许在证据不足时得到否定结论；这种可证伪性比预先承诺某个算法一定有效更符合研究报告的用途。

\hypertarget{ux7b2cux5341ux7ae0-ux7814ux7a76ux95eeux9898ux4e0eux7efcux5408ux7ed3ux8bba}{%
\section{第十章 研究问题与综合结论}\label{ux7b2cux5341ux7ae0-ux7814ux7a76ux95eeux9898ux4e0eux7efcux5408ux7ed3ux8bba}}

\hypertarget{ux9700ux8981ux4f18ux5148ux56deux7b54ux7684ux95eeux9898}{%
\subsection{10.1 需要优先回答的问题}\label{ux9700ux8981ux4f18ux5148ux56deux7b54ux7684ux95eeux9898}}

第一，如何建立同时反映地图高程、位置不确定性和空间相关误差的观测似然。第二，如何在平坦或重复地貌中保留合理的不确定性，而不是因大量相关测点产生虚假的高置信度。第三，如何使侧扫强度经过几何与辐射校正后提供独立约束，而非重复惯性拼接误差。

第四，如何把候选检索、鲁棒配准和连续多假设跟踪连接起来，使定位具有重新发现真实位置的能力。第五，如何将跨时间优化限定在已获得证据的候选内，避免错误绝对匹配污染历史轨迹。第六，如何以跨海域、跨航次及跨设备的独立数据评价学习表示和协方差的可靠性。

这些问题共同说明，海底地形导航不是一个独立``匹配模块''的精度竞赛，而是观测、地图、运动和关联的不确定性管理问题。方法的研究价值，应当由能否解决某一明确限制和能否提供可验证证据来判断。

\hypertarget{ux7efcux5408ux7ed3ux8bba}{%
\subsection{10.2 综合结论}\label{ux7efcux5408ux7ed3ux8bba}}

从经典地形导航到测深SLAM、概率地图和现代全局跟踪，研究主线始终是减少水下运动积分的漂移，并在环境信息不足时保持估计可信。已知地图的局部地形导航具有较成熟工程证据；测深和侧扫SLAM提供相对运动与地图一致性；学习型检索及全局多假设跟踪扩展了未知位置恢复能力，但跨区域鲁棒性仍需审慎评价。

多波束几何与侧扫强度融合具有明确研究潜力，其科学前提是正确的声学测量模型、时空标定、有效重叠和共享误差处理。地图分辨率、波束数量和算法复杂度均不能独立决定定位精度。平坦地形的不可观方向、重复地貌的多峰分布以及地图偏差的系统效应，需要被显式表达。

因此，合理研究路线应以物理观测和地图质量为基础，以局部可观性和全局歧义为两条分析轴，以概率可信度和失效恢复为系统评价核心。算法应按任务和地貌选用或组合，而不是将某一种陆地或水下方法视为无条件最优。这一结论来自多条公开研究路线的综合比较，不代表未经验证的性能结论。

\hypertarget{ux53c2ux8003ux6587ux732e}{%
\section{参考文献}\label{ux53c2ux8003ux6587ux732e}}

{[}1{]} Paull L, Saeedi S, Seto M, Li H. AUV Navigation and Localization: A Review. IEEE Journal of Oceanic Engineering, 2014, 39(1):131--149. \href{https://people.csail.mit.edu/lpaull/publications/Paull_JOE_2014.pdf}{公开来源}。

{[}2{]} Carreño S, Wilson P, Ridao P, Petillot Y. A survey on terrain based navigation for AUVs. OCEANS, 2010. DOI:10.1109/OCEANS.2010.5664372. \href{https://eprints.soton.ac.uk/153897/1/Carreno_2010_-_A_survey__on_TBN_for_AUVs.pdf}{公开来源}。

{[}3{]} Melo J, Matos A. Survey on advances on terrain based navigation for autonomous underwater vehicles. Ocean Engineering, 2017, 139:250--264. \href{https://doi.org/10.1016/j.oceaneng.2017.04.047}{公开来源}。

{[}4{]} Barkby S, Williams S B, Pizarro O, Jakuba M V. Bathymetric particle filter SLAM using trajectory maps. The International Journal of Robotics Research, 2012, 31(12):1409--1430. \href{https://doi.org/10.1177/0278364912459666}{公开来源}。

{[}5{]} Ånonsen K B, Hagen O K, Hegrenæs Ø, Hagen P E. The HUGIN AUV Terrain Navigation Module. OCEANS, 2013. \href{https://www.navlab.net/Publications/The_HUGIN_AUV_Terrain_Navigation_Module.pdf}{公开来源}。

{[}6{]} GEBCO. GEBCO global gridded bathymetry and source Type Identifier Grid. 官方全球海底地形数据说明. \href{https://www.gebco.net/data_and_products/gridded_bathymetry_data/}{公开来源}。

{[}7{]} NOAA NCEI. ETOPO Global Relief Model. 官方全球地形数据说明. \href{https://www.ncei.noaa.gov/products/etopo-global-relief-model}{公开来源}。

{[}8{]} GEBCO. Data set errata and known issues. 官方数据勘误. \href{https://www.gebco.net/data-products/gridded-bathymetry-data/data-set-errata}{公开来源}。

{[}9{]} Salavasidis G 等. Terrain-Aided Navigation With Coarse Maps---Toward an Arctic Crossing With an AUV. IEEE Journal of Oceanic Engineering, 2021, 46(4):1192--1212. \href{https://eprints.soton.ac.uk/480480/1/joe.pdf}{公开来源}。

{[}10{]} International Hydrographic Organization. IHO Standards for Hydrographic Surveys S-44 Edition 6.2.0. 国际海道测量标准, 2024版本. \href{https://iho.int/uploads/user/pubs/standards/s-44/S-44_Edition_6.2.0_adopted.pdf}{公开来源}。

{[}11{]} Torroba I, Sprague C I, Folkesson J. Fully-Probabilistic Terrain Modelling and Localization With Stochastic Variational Gaussian Process Maps. IEEE Robotics and Automation Letters, 2022, 7(4):8729--8736. \href{https://arxiv.org/abs/2203.10893}{公开来源}。

{[}12{]} Palomer A, Ridao P, Ribas D. Multibeam 3D Underwater SLAM with Probabilistic Registration. Sensors, 2016, 16(4):560. \href{https://pmc.ncbi.nlm.nih.gov/articles/PMC4851074/}{公开来源}。

{[}13{]} Lurton X, Lamarche G 主编 GeoHab Backscatter Working Group. Backscatter measurements by seafloor-mapping sonars Guidelines and Recommendations. 专业技术指南, 2015及后续资料. \href{https://geohab.org/backscatter-working-group/}{公开来源}。

{[}14{]} Zhang J, Xie Y, Ling L, Folkesson J. A Fully-automatic Side-scan Sonar SLAM Framework. 作者公开论文, 2023. 相关出版DOI:10.1049/rsn2.12500. \href{https://arxiv.org/abs/2304.01854}{公开来源}。

{[}15{]} Zhang J, Xie Y, Ling L, Folkesson J. A Dense Subframe-Based SLAM Framework With Side-Scan Sonar. IEEE Journal of Oceanic Engineering, 2025, 50(2):1087--1102. DOI:10.1109/JOE.2024.3503663. \href{https://arxiv.org/abs/2312.13802}{公开来源}。

{[}16{]} Aubard M, Madureira A, Teixeira L, Pinto J. Sonar-Based Deep Learning in Underwater Robotics: Overview, Robustness, and Challenges. IEEE Journal of Oceanic Engineering, 2025, 50(3):1866--1884. DOI:10.1109/JOE.2025.3531933. \href{https://arxiv.org/abs/2412.11840}{公开来源}。

{[}17{]} Pan X, Wu Y. Underwater Doppler Navigation With Self-calibration. The Journal of Navigation, 2016. DOI:10.1017/S0373463315000703. \href{https://arxiv.org/abs/1509.02054}{公开来源}。

{[}18{]} Xu Y, Zheng R, Zhang S, Liu M. Robust Inertial-aided Underwater Localization based on Imaging Sonar Keyframes. IEEE Transactions on Instrumentation and Measurement, 2022, 71:7501812. DOI:10.1109/TIM.2022.3156980. \href{https://arxiv.org/abs/2106.16032}{公开来源}。

{[}19{]} Zhang Q, Kim J. Terrain-Based Global Tracking Without Prior Pose Information Using Multibeam Echosounder Measurements. IEEE Transactions on Instrumentation and Measurement, 2026, 75:7506113. DOI:10.1109/TIM.2026.3716492. \href{https://doi.org/10.1109/TIM.2026.3716492}{出版入口}。本报告已核读用户提供的13页正式论文PDF。

{[}20{]} Segal A V, Haehnel D, Thrun S. Generalized-ICP. Robotics Science and Systems, 2009. \href{https://www.robots.ox.ac.uk/~avsegal/resources/papers/Generalized_ICP.pdf}{公开来源}。

{[}21{]} Barkby S, Williams S B, Pizarro O, Jakuba M V. Bathymetric SLAM with No Map Overlap using Gaussian Processes. IROS, 2011:1242--1248. \href{https://www-personal.acfr.usyd.edu.au/sbarkby/2011_iros_bpslam.pdf}{公开来源}。

{[}22{]} Yang H, Shi J, Carlone L. TEASER: Fast and Certifiable Point Cloud Registration. IEEE Transactions on Robotics, 2021, 37(2):314--333. \href{https://arxiv.org/abs/2001.07715}{公开来源}。

{[}23{]} Aoki K, Koide K, Oishi S, Yokozuka M, Banno A, Meguro J. 3D-BBS: Global Localization for 3D Point Cloud Scan Matching Using Branch-and-Bound Algorithm. ICRA, 2024; 作者版本arXiv:2310.10023. \href{https://arxiv.org/abs/2310.10023}{公开来源}。

{[}24{]} Ånonsen K B, Hallingstad O. Terrain Aided Underwater Navigation Using Point Mass and Particle Filters. IEEE/ION PLANS, 2006. DOI:10.1109/\allowbreak PLANS.\allowbreak 2006.\allowbreak 1650705. \href{https://www.navlab.net/Publications/Terrain_Aided_Underwater_Navigation_Using_Point_Mass_And_Particle_Filters.pdf}{公开来源}。

{[}25{]} Arulampalam M S, Maskell S, Gordon N, Clapp T. A tutorial on particle filters for online nonlinear/non-Gaussian Bayesian tracking. IEEE Transactions on Signal Processing, 2002, 50(2):174--188. \href{https://doi.org/10.1109/78.978374}{公开来源}。

{[}26{]} Fox D. KLD-Sampling: Adaptive Particle Filters. Advances in Neural Information Processing Systems 14, 2001. \href{https://papers.neurips.cc/paper_files/paper/2001/hash/c5b2cebf15b205503560c4e8e6d1ea78-Abstract.html}{公开来源}。

{[}27{]} Thrun S, Fox D, Burgard W, Dellaert F. Robust Monte Carlo localization for mobile robots. Artificial Intelligence, 2001, 128(1--2):99--141. \href{https://www.borg.cc.gatech.edu/projects/monte-carlo-localization.html}{公开来源}。

{[}28{]} Kaess M, Johannsson H, Roberts R, Ila V, Leonard J, Dellaert F. iSAM2: Incremental smoothing and mapping using the Bayes tree. The International Journal of Robotics Research, 2012, 31(2):216--235. \href{https://www.borg.cc.gatech.edu/papers/kaess12ijrr.html}{公开来源}。

{[}29{]} Olson E, Agarwal P. Inference on networks of mixtures for robust robot mapping. Robotics Science and Systems, 2012. \href{https://april.eecs.umich.edu/media/pdfs/olson2012rss.pdf}{公开来源}。

{[}30{]} Xie Y, Zhang J, Bore N, Folkesson J. NeuRSS: Enhancing AUV Localization and Bathymetric Mapping with Neural Rendering for Sidescan SLAM. 作者公开论文, 2024, arXiv:2405.05807. \href{https://arxiv.org/abs/2405.05807}{公开来源}。

{[}31{]} Xu W, Cai Y, He D, Lin J, Zhang F. FAST-LIO2: Fast Direct LiDAR-Inertial Odometry. IEEE Transactions on Robotics, 2022, 38(4):2053--2073. \href{https://arxiv.org/abs/2107.06829}{公开来源}。

{[}32{]} Zheng C 等. FAST-LIVO2: Fast, Direct LiDAR--Inertial--Visual Odometry. IEEE Transactions on Robotics, 2025, 41:326--346. \href{https://arxiv.org/abs/2408.14035}{公开来源}。

{[}33{]} MBARI. Terrain Relative Navigation. 官方技术项目资料. \href{https://www.mbari.org/project/terrain-relative-navigation/}{公开来源}。

{[}34{]} Caress D W 等. MB-System terrain-relative navigation source code. 官方开源软件资料. \href{https://github.com/dwcaress/MB-System}{公开来源}。

{[}35{]} Tan J, Torroba I, Xie Y, Folkesson J. Data-driven Loop Closure Detection in Bathymetric Point Clouds for Underwater SLAM. 作者公开论文, 2022, arXiv:2209.08578; 相关ICRA 2023出版. \href{https://arxiv.org/abs/2209.08578}{公开来源}。

{[}36{]} Torroba I, Sprague C I, Bore N, Folkesson J. PointNetKL: Deep Inference for GICP Covariance Estimation in Bathymetric SLAM. IEEE Robotics and Automation Letters, 2020, 5(3):4078--4085; 作者版本arXiv:2003.10931. \href{https://arxiv.org/abs/2003.10931}{公开来源}。

{[}37{]} Uy M A, Lee G H. PointNetVLAD: Deep Point Cloud Based Retrieval for Large-Scale Place Recognition. CVPR, 2018. \href{https://arxiv.org/abs/1804.03492}{公开来源}。

{[}38{]} Komorowski J. MinkLoc3D: Point Cloud Based Large-Scale Place Recognition. WACV, 2021. \href{https://arxiv.org/abs/2011.04530}{公开来源}。

{[}39{]} Roberts D R 等. Cross-validation strategies for data with temporal, spatial, hierarchical, or phylogenetic structure. Ecography, 2017, 40:913--929. \href{https://doi.org/10.1111/ecog.02881}{公开来源}。

{[}40{]} Harris P T, Macmillan-Lawler M, Rupp J, Baker E K. Geomorphology of the oceans. Marine Geology, 2014, 352:4--24. DOI:10.1016/j.margeo.2014.01.011. \href{https://bluehabitats.org/}{公开来源}。

{[}41{]} Ma D, Ma T, Li Y, Zhang Q, Ling Y, Liao Y. Underwater multizonotope terrain-aided navigation method with coarse map based on set-membership filter. Journal of Field Robotics, 2025, 42(1):242--256. \href{https://onlinelibrary.wiley.com/doi/abs/10.1002/rob.22407}{公开来源}。

{[}42{]} Torroba I, Cella M, Terán A, Rolleberg N, Folkesson J. Online Stochastic Variational Gaussian Process Mapping for Large-Scale Bathymetric SLAM in Real Time. IEEE Robotics and Automation Letters, 2023. DOI:10.1109/LRA.2023.3264750. \href{https://arxiv.org/abs/2211.05601}{公开来源}。

{[}43{]} Xie Y, Bore N, Folkesson J. Bathymetric Reconstruction From Sidescan Sonar With Deep Neural Networks. IEEE Journal of Oceanic Engineering, 2023, 48(2):372--383. DOI:10.1109/JOE.2022.3220330. \href{https://arxiv.org/abs/2206.07810}{公开来源}。

{[}44{]} Xie Y, Bore N, Folkesson J. Neural Network Normal Estimation and Bathymetry Reconstruction from Sidescan Sonar. 作者公开论文, 2022, arXiv:2206.07819. \href{https://arxiv.org/abs/2206.07819}{公开来源}。

{[}45{]} Xie Y, Bore N, Folkesson J. Towards Differentiable Rendering for Sidescan Sonar Imagery. IEEE/OES AUV Symposium, 2022. DOI:10.1109/\allowbreak AUV53081.\allowbreak 2022.\allowbreak 9965917. \href{https://arxiv.org/abs/2206.07821}{公开来源}。

{[}46{]} NOAA NCEI. Coastal Elevation Models and Continuously Updated Digital Elevation Models. 官方沿海地形与水深数据产品说明. \href{https://www.ncei.noaa.gov/products/coastal-elevation-models}{公开来源}。

{[}47{]} Global Multi-Resolution Topography synthesis. GMRT Overview and Technical Details. 官方多分辨率海底地形产品说明. \href{https://gmrt.org/about/index.php}{公开来源}。

{[}48{]} EMODnet. EMODnet Bathymetry. 官方欧洲海域DTM、高分辨率产品与质量层说明. \href{https://emodnet.ec.europa.eu/en/bathymetry}{公开来源}。

{[}49{]} AusSeabed. Bathymetry data. 官方澳大利亚海底数据与测量覆盖说明. \href{https://www.ausseabed.gov.au/data/bathymetry}{公开来源}。

{[}50{]} Zhang Q, Kim J. TTT SLAM: A feature-based bathymetric SLAM framework. Ocean Engineering, 2024, 294:116777. DOI:10.1016/j.oceaneng.2024.116777. \href{https://www.sciencedirect.com/science/article/pii/S0029801824001148}{公开来源}。

{[}51{]} Yang H, Antonante P, Tzoumas V, Carlone L. Graduated Non-Convexity for Robust Spatial Perception: From Non-Minimal Solvers to Global Outlier Rejection. IEEE Robotics and Automation Letters, 2020. DOI:10.1109/LRA.2020.2965893. \href{https://arxiv.org/abs/1909.08605}{公开来源}。

{[}52{]} Cadena C, Carlone L, Carrillo H, Latif Y, Scaramuzza D, Neira J, Reid I, Leonard J J. Past, Present, and Future of Simultaneous Localization and Mapping: Towards the Robust-Perception Age. IEEE Transactions on Robotics, 2016, 32(6):1309--1332. DOI:10.1109/TRO.2016.2624754. \href{https://arxiv.org/abs/1606.05830}{公开来源}。

\hypertarget{ux6587ux732eux8bc1ux636eux4e0eux9605ux8bfbux8303ux56f4}{%
\section{文献证据与阅读范围}\label{ux6587ux732eux8bc1ux636eux4e0eux9605ux8bfbux8303ux56f4}}

以下区分公开全文、可见方法章节、摘要和软件资料。摘要层面的文献只用于陈述其公开问题定义、总体方法和证据类型，不据此补写未公开公式、参数或逐项实验结果。通用视觉和陆地研究仅作为方法来源，不等同于水下有效性证明。

\begin{longtable}[]{@{}p{0.07\linewidth}p{0.29\linewidth}p{0.58\linewidth}@{}}
\toprule
编号 & 证据类型 & 本报告使用范围\tabularnewline
\midrule
\endhead
1 & 水下导航综述 & 作者公开全文\tabularnewline
2 & 水下地形导航综述 & 作者公开全文\tabularnewline
3 & 水下地形导航综述 & 出版摘要与机构公开全文\tabularnewline
4 & 测深SLAM直接研究 & 公开摘要及作者相关方法论文\tabularnewline
5 & 工程系统与海试证据 & 作者公开全文\tabularnewline
6 & 地图来源与质量 & 官方数据说明 历史产品与新版应分别引用\tabularnewline
7 & 地图来源与质量 & 官方数据说明\tabularnewline
8 & 地图质量边界 & 官方勘误说明\tabularnewline
9 & 粗地图导航直接研究 & 作者公开全文\tabularnewline
10 & 测量不确定度术语 & 官方标准全文\tabularnewline
11 & 概率海底地图直接研究 & 作者公开全文与官方实现\tabularnewline
12 & 多波束概率配准直接研究 & 开放全文\tabularnewline
13 & 声学强度物理基础 & 官方指南与相关章节\tabularnewline
14 & 侧扫SLAM直接研究 & 作者公开全文与官方实现\tabularnewline
15 & 侧扫SLAM直接研究 & 开放全文与官方实现\tabularnewline
16 & 声呐学习综述 & 开放作者全文\tabularnewline
17 & DVL标定与可观性直接研究 & 作者公开全文\tabularnewline
18 & 成像声呐惯性融合直接研究 & 作者公开全文与相关出版DOI核对\tabularnewline
19 & 全局地形跟踪直接研究 & 用户提供IEEE正式论文PDF，13页全文核读；原方法及实验边界已核对\tabularnewline
20 & 概率配准基础 & 作者公开全文\tabularnewline
21 & 测深SLAM直接研究 & 作者公开全文\tabularnewline
22 & 鲁棒全局配准基础 & 作者公开全文\tabularnewline
23 & 全局搜索参考 & 作者摘要与官方实现说明\tabularnewline
24 & 水下地形导航直接研究 & 作者公开全文\tabularnewline
25 & 粒子滤波理论基础 & 公开正文与正式出版入口\tabularnewline
26 & 粒子数自适应基础 & 官方论文摘要与全文入口\tabularnewline
27 & 全局定位概率基础 & 作者实验室方法与书目说明\tabularnewline
28 & 增量因子图基础 & 作者实验室书目与方法说明\tabularnewline
29 & 鲁棒多假设图优化基础 & 作者公开全文\tabularnewline
30 & 侧扫重建与SLAM直接研究 & 作者公开全文\tabularnewline
31 & 陆地方法迁移参考 & 作者公开全文与官方实现\tabularnewline
32 & 陆地多模态方法迁移参考 & 作者公开全文与官方实现\tabularnewline
33 & 工程系统证据 & 官方项目说明\tabularnewline
34 & 可复现软件证据 & 官方仓库与选定TRN源文件\tabularnewline
35 & 测深地点识别直接研究 & 作者公开全文与作者书目核对\tabularnewline
36 & 测深配准不确定性直接研究 & 作者公开全文与机构出版核对\tabularnewline
37 & 陆地点云检索参考 & 作者摘要与公开论文入口\tabularnewline
38 & 陆地点云检索参考 & 作者摘要与公开论文入口\tabularnewline
39 & 空间独立评价基础 & 正式出版方法综述\tabularnewline
40 & 全球海底地貌分类基础 & 作者机构合作项目与NOAA地貌图层说明 不是论文全文核读\tabularnewline
41 & 粗地图集合估计直接研究 & 出版摘要与期刊信息\tabularnewline
42 & 在线测深概率建图直接研究 & 作者公开论文摘要与官方实现说明\tabularnewline
43 & 侧扫几何重建直接研究 & 作者公开全文与机构出版核对\tabularnewline
44 & 声学表面法向重建研究 & 作者公开全文\tabularnewline
45 & 声学可微模型研究 & 作者公开全文\tabularnewline
46 & 真实海底地形数据来源 & 官方产品说明 具体数据精度需逐项元数据核查\tabularnewline
47 & 真实海底地形数据来源 & 官方产品与处理说明 不把全球底图视作全域高分测深\tabularnewline
48 & 真实海底地形数据来源 & 官方产品说明 高分辨率限特定区域 原始测量获取条件另核查\tabularnewline
49 & 真实海底地形数据来源 & 官方数据门户说明 不推断每个航次均含开放强度\tabularnewline
50 & 测深特征与鲁棒SLAM直接研究 & 出版摘要、引言与可见方法片段 不声称全文复现\tabularnewline
51 & 鲁棒估计方法基础 & 作者摘要与公开版本入口 不据此声称无条件全局最优\tabularnewline
52 & 机器人SLAM综述 & 作者公开版本入口及摘要 用于SLAM问题分类\tabularnewline
\bottomrule
\end{longtable}
